\subsection{NODAL VECTORS AND INVERSE ROOTS}

Lemma 2. Let S be a closed subgroup of T and Z(S) its normalizer in G.

\begin{enumerate}
\def\labelenumi{(\roman{enumi})}
\item
  \(\mathrm{R}(\mathrm{Z}(\mathrm{S})_0,\mathrm{T})\) is the set of \(\alpha \in \mathrm{R}(\mathrm{G},\mathrm{T})\) such that \(\alpha (\mathrm{S}) = \{1\}\) .
\item
  The centre of \(\mathrm{Z(S)_0}\) is the intersection of the \(\operatorname {Ker}\alpha\) for \(\alpha \in \mathrm{R}(\mathrm{Z(S)}_0,\mathrm{T})\) .
\item
  If S is connected, Z(S) is connected.
\end{enumerate}

The Lie algebra \(\mathrm{L}(\mathrm{Z}(\mathrm{S}))_{(\mathbf{C})}\) consists of the invariants of S on \(g_{C}\) (Chap. III, §9, no. 3, Prop. 8), and hence is the direct sum of \(t_{C}\) and the \(g^{\alpha}\) for which \(\alpha(\mathrm{S})=\{1\}\) , hence (i). Assertion (ii) follows from Prop. 8 (no. 4), and assertion (iii) has already been proved (§2, no. 2, Cor. 5 of Th. 2).

THEOREM 1. Let \(\alpha \in \mathrm{R}(\mathrm{G},\mathrm{T})\) . The centralizer \(Z_{\alpha}\) of the kernel of \(\alpha\) is a connected closed subgroup of \(\mathrm{G}\) ; its centre is \(\operatorname{Ker} \alpha\) ; its derived group \(\mathrm{D}(Z_{\alpha}) = S_{\alpha}\) is a connected closed semi-simple subgroup of \(\mathrm{G}\) of rank 1. We have \(\mathrm{R}(Z_{\alpha},\mathrm{T}) = \{\alpha, -\alpha\}\) and \(\dim Z_{\alpha} = \dim T + 2\) .

Let \(Z_{\alpha}^{\prime}\) be the centralizer of \((\mathrm{Ker}\alpha)_0\) . By Lemma 2, this is a connected closed subgroup of G, and \(\mathrm{R}(Z_{\alpha}^{\prime},\mathrm{T})\) is the set of \(\beta \in \mathrm{R}(\mathrm{G},\mathrm{T})\) such that \(\beta ((\mathrm{Ker}\alpha)_0) = \{1\}\) . Clearly, \(\{\alpha , - \alpha \} \subset \mathrm{R}(Z_{\alpha}^{\prime},\mathrm{T})\) . Conversely, let \(\beta \in \mathrm{R}(Z_{\alpha}^{\prime},\mathrm{T})\) ; since \((\mathrm{Ker}\alpha)_0\) is of finite index in \(\mathrm{Ker}\alpha\) , there exists an integer \(r\neq 0\) such that \(t^{r\beta} = 1\) for \(t\in \mathrm{Ker}\alpha\) . From the exactness of the sequence

\[
0 \longrightarrow \mathbf {Z} \longrightarrow \mathrm{X(T)} \longrightarrow \mathrm{X(Ker} \alpha) \longrightarrow 0
\]

corresponding by duality to the exact sequence

\[
0 \longrightarrow \operatorname{Ker} \alpha \longrightarrow \mathrm{T} \stackrel {{\alpha}} {{\longrightarrow}} \mathrm{U} \longrightarrow 0,
\]

it follows that \(r\beta\) is a multiple of \(\alpha\) ; by Chap. VIII, §2, no. 2, Th. 2 (i), this implies that \(\beta \in \{\alpha, -\alpha\}\) . Thus, \(\mathrm{R}(Z_{\alpha}', T) = \{\alpha, -\alpha\}\) . It follows (Lemma 2) that the centre of \(Z_{\alpha}'\) is \(\operatorname{Ker} \alpha\) , so \(Z_{\alpha}' = Z_{\alpha}\) . Finally, by Cor. 1 of Prop. 4 (§1, no. 4), \(\mathrm{D}(\mathrm{Z}_{\alpha})\) is a connected closed semi-simple subgroup of \(\mathrm{G}\) ; it is of rank 1 because \(\mathcal{DL}(\mathrm{Z}_{\alpha})_{(\mathbf{C})} = \mathfrak{g}^{\alpha} + \mathfrak{g}^{-\alpha} + [\mathfrak{g}^{\alpha},\mathfrak{g}^{-\alpha}]\) .

COROLLARY. There exists a morphism of Lie groups \(\nu : \mathbf{SU}(2, \mathbf{C}) \to \mathrm{G}\) with the following properties:

\begin{enumerate}
\def\labelenumi{\alph{enumi})}
\item
  The image of \(\nu\) commutes with the kernel of \(\alpha\) .
\item
  For all \(a \in U\) , we have \(\nu\left(\begin{array}{cc}a & 0 \\ 0 & \bar{a}\end{array}\right) \in T\) and \(\alpha \circ \nu\left(\begin{array}{cc}a & 0 \\ 0 & \bar{a}\end{array}\right) = a^{2}\) .
\end{enumerate}

If \(\nu_{1}\) and \(\nu_{2}\) are two morphisms from \(\mathbf{SU}(2,\mathbf{C})\) to \(\mathrm{G}\) with the preceding properties, there exists \(a\in \mathbf{U}\) such that \(\nu_{2} = \nu_{1}\circ \operatorname {Int}\left( \begin{array}{cc}a & 0\\ 0 & \bar{a} \end{array} \right)\) .

By Th. 1 and Prop. 6 of §3, no. 6, there exists a morphism of Lie groups \(\nu : \mathbf{SU}(2, \mathbf{C}) \to \mathrm{S}_{\alpha}\) that is surjective with discrete kernel. Then \(\nu^{-1}(\mathrm{T} \cap \mathrm{S}_{\alpha})\) is a maximal torus of \(\mathbf{SU}(2, \mathbf{C})\) (§2, no. 3, Prop. 1). Since the maximal tori of \(\mathbf{SU}(2, \mathbf{C})\) are conjugate (§2, no. 2, Th. 2), we can assume, replacing \(\nu\) by \(\nu \circ \operatorname{Int}s\) (with \(s \in \mathbf{SU}(2, \mathbf{C})\) ) if necessary, that \(\nu^{-1}(\mathrm{T} \cap \mathrm{S}_{\alpha})\) is the group of diagonal matrices in \(\mathbf{SU}(2, \mathbf{C})\) . Then \(\nu\begin{pmatrix} a & 0 \\ 0 & \bar{a} \end{pmatrix} \in \mathrm{T}\) for all \(a \in \mathbf{U}\) , and the map

\[
\left( \begin{array}{c c} a & 0 \\ 0 & \bar {a} \end{array} \right) \mapsto \alpha \circ \nu \left( \begin{array}{c c} a & 0 \\ 0 & \bar {a} \end{array} \right)
\]

is a root of \(\mathbf{SU}(2,\mathbf{C})\) , and hence is equal to one of the two maps \(\begin{pmatrix}a&0\\ 0&\bar{a}\end{pmatrix}\mapsto a^{2}\) or \(\begin{pmatrix}a&0\\ 0&\bar{a}\end{pmatrix}\mapsto a^{-2}\) (§3, no. 6, formulas (19)). In the first case, the homomorphism \(\nu\) has the required properties; in the second case, the homomorphism \(\nu\circ Int\theta\) has them (loc. cit., formulas (18)).

If \(\nu_{1}\) and \(\nu_{2}\) are two morphisms from \(\mathbf{SU}(2,\mathbf{C})\) to G satisfying the stated conditions, they both map \(\mathbf{SU}(2,\mathbf{C})\) into \(S_{\alpha}\) (condition a)), hence are both universal coverings of \(S_{\alpha}\) . Hence, there exists an automorphism \(\varphi\) of \(\mathbf{SU}(2,\mathbf{C})\) such that \(\nu_{2} = \nu_{1} \circ \varphi\) , and we conclude by using Prop. 9 of no. 4.

It follows from the preceding corollary that the homomorphism \(\nu_{\mathrm{T}}\) from \(\mathbf{U}\) to \(\mathrm{T}\) , defined by \(\nu_{\mathrm{T}}(a) = \nu \left( \begin{array}{cc}a & 0\\ 0 & \bar{a} \end{array} \right)\) for \(a\in \mathbf{U}\) , is independent of the choice of \(\nu\) . Denote by \(K_{\alpha}\in \Gamma (\mathrm{T})\) the image under \(\Gamma (\nu_{\mathrm{T}})\) of the element \(2\pi i\) of \(\Gamma (\mathbf{U}) = 2\pi i\mathbf{Z}\) ; it is called the nodal vector associated to the root \(\alpha\) . We have \(\langle \alpha ,K_{\alpha}\rangle = 2\) , in other words (no. 2, formula (2)) \(\delta (\alpha)(K_{\alpha}) = 4\pi i\) ; since \(K_{\alpha}\) belongs to the intersection of \(\mathfrak{t}\) and the \(\mathrm{L}(\mathrm{S}_{\alpha})_{(\mathbf{C})}\) , we have

\[
K _ {\alpha} = 2 \pi i H _ {\delta (\alpha)},\tag{13}
\]

where \(H_{\delta(\alpha)}\) is the inverse root associated to the root \(\delta(\alpha)\) of \((\mathfrak{g}_{\mathbf{C}}, \mathfrak{t}_{\mathbf{C}})\) (Chap. VIII, §2, no. 2). In other words, when \(\Gamma(\mathrm{T}) \otimes \mathbf{R}\) is identified with the dual of \(\mathrm{X(T)}\otimes \mathbf{R}\) via the pairing \(\langle ,\rangle ,K_{\alpha}\) is identified with the inverse root \(\alpha^{\vee}\in (\mathrm{X(T)}\otimes \mathbf{R})^{*}\) .

Remark. For all \(x \in R\) , we have

\[
\nu \left( \begin{array}{c c} \exp (2 \pi i x) & 0 \\ 0 & \exp (- 2 \pi i x) \end{array} \right) = \nu_ {\mathrm{T}} (e ^ {2 \pi i x}) = \exp (x K _ {\alpha}).\tag{14}
\]

In particular:

\[
\nu \left( \begin{array}{c c} - 1 & 0 \\ 0 & - 1 \end{array} \right) = \nu_ {\mathrm{T}} (- 1) = \exp \Bigl (\frac {1}{2} K _ {\alpha} \Bigr).\tag{15}
\]

It follows that \(\nu\) is injective if and only if \(K_{\alpha} \notin 2\Gamma(T)\) , in other words if there exists \(\lambda \in \mathrm{X}(T)\) such that \(\langle \lambda, K_{\alpha} \rangle \notin 2Z\) . When \(g_{C}\) is simple, \(\nu\) is injective unless \(g_{C}\) is of type \(B_{n}\) , \(C(G) = \{1\}\) and \(\alpha\) is a short root (cf.~Chap. VI, Plates).

In the remainder of this paragraph we denote by \(\mathrm{R}^{\vee}(\mathrm{G},\mathrm{T})\) the set of nodal vectors \(K_{\alpha}\) for \(\alpha\in\mathrm{R}(\mathrm{G},\mathrm{T})\) . This is a subset of \(\varGamma(\mathrm{T})\) that the canonical injection of \(\varGamma(\mathrm{T})\) into \(t_{C}\) identifies with the homothety with ratio \(2\pi i\) of the inverse root system \(\mathrm{R}^{\vee}(\mathfrak{g}_{\mathbf{C}},\mathfrak{t}_{\mathbf{C}})=\{H_{\delta(\alpha)}\}\) of \(\delta(\mathrm{R})\) . It follows that \(\mathrm{R}^{\vee}(\mathrm{G},\mathrm{T})\) generates the R-vector space \(\mathrm{L}(\mathrm{T}\cap\mathrm{D}(\mathrm{G}))\) , and hence that its orthogonal complement in \(\mathrm{X}(\mathrm{T})\) is \(\mathrm{X}(\mathrm{T}/(\mathrm{T}\cap\mathrm{D}(\mathrm{G})))\) .

Denote by \(\mathrm{Aut(T)}\) the group of automorphisms of the Lie group T; the Weyl group \(W = W_{G}(T)\) (§2, no. 5) can be identified with a subgroup of \(\mathrm{Aut(T)}\) . On the other hand, recall (Chap. VIII, §2, no. 2, Remark 4) that the Weyl group \(W(\mathfrak{g}_{\mathbf{C}}, t_{\mathbf{C}})\) of the split reductive algebra ( \(\mathfrak{g}_{\mathbf{C}}, t_{\mathbf{C}}\) ) operates on \(t_{\mathbf{C}}\) , and thus is canonically identified with a subgroup of \(\mathbf{GL}(t_{\mathbf{C}})\) .

PROPOSITION 10. The map \(u \mapsto \mathrm{L}(u)_{(\mathbf{C})}\) from \(\operatorname{Aut}(\mathrm{T})\) to \(\mathbf{GL}(\mathfrak{t}_{\mathbf{C}})\) induces an isomorphism from \(\mathrm{W}\) to the Weyl group of the split reductive Lie algebra \((\mathfrak{g}_{\mathbf{C}}, \mathfrak{t}_{\mathbf{C}})\) . For all \(\alpha \in \mathrm{R}\) , \(\mathrm{W}_{\mathrm{Z}_{\alpha}}(\mathrm{T})\) is of order 2, and the image under the preceding isomorphism of the non-identity element of \(\mathrm{W}_{\mathrm{Z}_{\alpha}}(\mathrm{T})\) is the reflection \(s_{H_{\delta(\alpha)}}\) .

The map under consideration is injective. It remains to show that its image is equal to \(\mathrm{W}(\mathfrak{g}_{\mathbf{C}}, t_{\mathbf{C}})\) .

Let \(g \in \mathrm{N}_{\mathrm{G}}(\mathrm{T})\) . With the notations in Chap. VIII, §5, no. 2, we have \(\operatorname{Ad} g \in \operatorname{Aut}(\mathfrak{g}_{\mathbf{C}}, \mathfrak{t}_{\mathbf{C}}) \cap \operatorname{Int}(\mathfrak{g}_{\mathbf{C}})\) , so \(\operatorname{Ad} g \in \operatorname{Aut}_{0}(\mathfrak{g}_{\mathbf{C}}, \mathfrak{t}_{\mathbf{C}})\) (loc. cit., no. 5, Prop. 11). By loc. cit., no. 2, Prop. 4, the automorphism of \(\mathfrak{t}_{\mathbf{C}}\) induced by \(\operatorname{Ad} g\) belongs to \(\mathrm{W}(\mathfrak{g}_{\mathbf{C}}, \mathfrak{t}_{\mathbf{C}})\) . Thus, the image of \(\mathrm{W}\) in \(\mathbf{GL}(\mathfrak{t}_{\mathbf{C}})\) is contained in \(\mathrm{W}(\mathfrak{g}_{\mathbf{C}}, \mathfrak{t}_{\mathbf{C}})\) .

Let \(\alpha \in \mathrm{R}(\mathrm{G},\mathrm{T})\) , and let \(\nu : \mathbf{SU}(2,\mathbf{C}) \to \mathrm{G}\) be a morphism of Lie groups having the properties in the Cor. of Th. 1. The image under \(\nu\) of the element \(\theta\) of \(\mathbf{SU}(2,\mathbf{C})\) has the following properties (§3, no. 6, formulas (17)):

\[
a) (\operatorname{Int} \nu (\theta)) (t) = t \quad \text { if } t \in \operatorname{Ker} \alpha ,
\]

\[
b) (\operatorname{Int} \nu (\theta)) (t) = t ^ {- 1} \quad \text { if } t \in \mathrm{T} \cap \mathrm{S} _ {\alpha}.
\]

It follows that \(\operatorname{Ad}\nu(\theta)\) induces the identity on \(\operatorname{Ker}\delta(\alpha)\subset\mathfrak{t}_{\mathbf{C}}\) , and induces the map \(x\mapsto-x\) on \([\mathfrak{g}^{\alpha},\mathfrak{g}^{-\alpha}]\) , hence coincides with the reflection \(s_{H_{\delta(\alpha)}}\) . Thus, the image of W contains all the \(s_{H_{\delta(\alpha)}}\) , and hence is equal to \(\mathrm{W}(\mathfrak{g}_{\mathbf{C}}, t_{\mathbf{C}})\) . In particular \(\mathrm{W}_{\mathrm{Z}_{\alpha}}(\mathrm{T})\) is of order 2, and hence consists of the identity and \(\operatorname{Int}\nu(\theta)\) . This completes the proof of the proposition.

COROLLARY. Assume that G is semi-simple. Then every element of G is the commutator of two elements of G.

Let \(c\) be a Coxeter transformation of the Weyl group \(\mathrm{W}(\mathfrak{g}_{\mathbf{C}},\mathfrak{t}_{\mathbf{C}})\) (Chap. V, §6, no. 1), and let \(n\) be an element of \(\mathrm{N}_{\mathrm{G}}(\mathrm{T})\) whose class in \(\mathrm{W}\) is identified with \(c\) by the isomorphism defined in the proposition. Denote by \(f_{c}\) the morphism \(t\mapsto (n,t)\) from \(\mathrm{T}\) to \(\mathrm{T}\) ; for \(x\in \mathfrak{t}_{\mathbf{C}}\) , we have \(\mathrm{L}(f_c)_{(\mathbf{C})}(x) = (\mathrm{Ad}n)(x) - x = c(x) - x\) .

By Th. 1 of Chap. V, §6, no. 2, the endomorphism \(c\) of \(\mathbf{t}_{\mathbf{C}}\) has no eigenvalue equal to 1. Consequently, \(\mathrm{L}(f_c)\) is surjective, and hence so is \(f_c\) . It follows that every element of T is the commutator of two elements of G, which implies the corollary in view of Th. 2, §2, no. 2.

\subsection{FUNDAMENTAL GROUP}

In the following proposition, \(f(\mathrm{G},\mathrm{T})\) denotes the homomorphism from \(\varGamma(\mathrm{T})\) to \(\pi_{1}(\mathrm{G})\) that is the composite of the canonical isomorphism from \(\varGamma(\mathrm{T})\) to \(\pi_{1}(\mathrm{T})\) (no. 2, Remark 3) and the homomorphism \(\pi_{1}(\iota)\) , where \(\iota\) is the canonical injection \(T\to G\) .

PROPOSITION 11. The homomorphism \(f(\mathrm{G},\mathrm{T}) : \Gamma(\mathrm{T}) \to \pi_1(\mathrm{G})\) is surjective. Its kernel is the subgroup \(\mathrm{N}(\mathrm{G},\mathrm{T})\) of \(\Gamma(\mathrm{T})\) generated by the family of nodal vectors \((K_{\alpha})_{\alpha \in \mathrm{R}(\mathrm{G},\mathrm{T})}\) .

The homomorphism \(f(\mathrm{G},\mathrm{T})\) is surjective by Prop. 3 ( \(\S2\) , no. 4). We denote by \(A(\mathrm{G},\mathrm{T})\) the assertion: ``the kernel of \(f(\mathrm{G},\mathrm{T})\) is generated by the \(K_{\alpha}\) '' which it remains to prove, and distinguish several cases:

\begin{enumerate}
\def\labelenumi{\alph{enumi})}
\item
  G is simply-connected. Let \(\rho : \mathfrak{g}_{\mathbf{C}} \to \mathfrak{gl}(\mathrm{V})\) be a linear representation of \(\mathfrak{g}_{\mathbf{C}}\) on a finite dimensional complex vector space V. Restricting to \(\mathfrak{g}\) , we obtain a representation of \(\mathfrak{g}\) on the real vector space \(V_{(\mathbf{R})}\) ; since G is simply-connected, there exists an analytic linear representation \(\pi\) of G on \(V_{(\mathbf{R})}\) such that \(\rho = L(\pi)\) . It follows from Prop. 7 of no. 3 that the image \(\delta(X(T))\) of X(T) in \(\mathfrak{t}_{\mathbf{C}}^{*}\) contains all the weights of \(\rho\) on V. This being true for every representation \(\rho\) of \(\mathfrak{g}_{\mathbf{C}}\) , it follows from Chap. VIII, §7, no. 2, Th. 1 that \(\delta(X(T))\) contains the group of weights of \(\delta(R)\) , which is by definition the set of \(\lambda \in \mathfrak{t}_{\mathbf{C}}^{*}\) such that \(\lambda(H_{\delta(\alpha)}) \in \mathbf{Z}\) for all \(\alpha \in R\) , in other words, \(\lambda(K_{\alpha}) \in 2\pi i\mathbf{Z}\) for all \(\alpha \in R\) . Thus, the group X(T) contains all the elements \(\lambda\) of X(T) \(\otimes\) Q such that \(\langle\lambda, K_{\alpha} \rangle \in \mathbf{Z}\) for all \(\alpha \in R\) , which implies by duality that \(\Gamma(T)\) is generated by the \(K_{\alpha}\) , hence the assertion \(A(G,T)\) .
\item
  G is the direct product of a simply-connected group \(G'\) and a torus S. Then T is the direct product of a maximal torus \(T'\) of \(G'\) with S, \(\varGamma(\mathrm{T})\) can be identified with \(\varGamma(\mathrm{T}')\times\varGamma(\mathrm{S})\) , \(\pi_{1}(\mathrm{G})\) with \(\pi_{1}(\mathrm{G}')\times\pi_{1}(\mathrm{S})\) , and \(f(\mathrm{G},\mathrm{T})\) with the homomorphisms with components \(f(\mathrm{G}^{\prime},\mathrm{T}^{\prime})\) and \(f(\mathrm{S},\mathrm{S})\) . Since \(f(\mathrm{S},\mathrm{S})\) is bijective, the canonical map \(\varGamma(\mathrm{T}^{\prime})\to\varGamma(\mathrm{T})\) maps \(\operatorname{Ker}f(\mathrm{G}^{\prime},\mathrm{T}^{\prime})\) bijectively onto \(\operatorname{Ker}f(\mathrm{G},\mathrm{T})\) . Moreover, the \(K_{\alpha}\) belong to the Lie algebra of the derived group \(G^{\prime}\) of G, hence to the image of \(\varGamma(\mathrm{T}^{\prime})\) , so it is immediate that \(A(\mathrm{G}^{\prime},\mathrm{T}^{\prime})\) implies \(A(\mathrm{G},\mathrm{T})\) , hence assertion \(A(\mathrm{G},\mathrm{T})\) , in view of a).
\item
  General case. There exists a surjective morphism \(p : G' \to G\) with finite kernel, where \(G'\) is the direct product of a simply-connected group by a torus ( \(\S1\) , no. 4, Prop. 4). If \(T'\) is the inverse image of T in \(G'\) (this is a maximal torus of \(G'\) by \(\S2\) , no. 3, Prop. 1), and N the kernel of p, we have exact sequences \(0 \to N \to G' \to G \to 0\) and \(0 \to N \to T' \to T \to 0\) , hence a commutative diagram with exact rows (no. 2, Remark 1 and General Topology, Chap. XI, in preparation)
\end{enumerate}

\[
\begin{array}{c c c c c c c c c} 0 & \longrightarrow & \Gamma (\mathrm{T} ^ {\prime}) & \longrightarrow & \Gamma (\mathrm{T}) & \longrightarrow & \mathrm{N} & \longrightarrow & 0 \\ & & \Big \downarrow f (\mathrm{G} ^ {\prime}, \mathrm{T} ^ {\prime}) & & \Big \downarrow f (\mathrm{G}, \mathrm{T}) & & \Big \downarrow \mathrm{Id} _ {\mathrm{N}} \\ 0 & \longrightarrow & \pi_ {1} (\mathrm{G} ^ {\prime}) & \longrightarrow & \pi_ {1} (\mathrm{G}) & \longrightarrow & \mathrm{N} & \longrightarrow & 0. \end{array}
\]

It follows immediately from the snake diagram (Algebra, Chap. X, p.~4, Prop. 2) that \(A(\mathrm{G}', \mathrm{T}')\) implies \(A(\mathrm{G}, \mathrm{T})\) , hence the proposition, in view of \(b\) ).

COROLLARY 1. G is simply-connected if and only if the family \((K_{\alpha})_{\alpha \in \mathrm{R}(\mathrm{G},\mathrm{T})}\) generates \(\Gamma (\mathrm{T})\) .

COROLLARY 2. Let H be a connected closed subgroup of G containing T; there is an exact sequence

\[
0 \longrightarrow \mathrm{N} (\mathrm{H}, \mathrm{T}) \longrightarrow \mathrm{N} (\mathrm{G}, \mathrm{T}) \longrightarrow \pi_ {1} (\mathrm{H}) \longrightarrow \pi_ {1} (\mathrm{G}) \longrightarrow 0.
\]

This follows from Algebra, Chap. X, p.~4, Prop. 2 (snake diagram), applied to the commutative diagram

\[
\begin{array}{c c c c c c c c c} 0 & \longrightarrow & \mathrm{N(H,T)} & \longrightarrow & \varGamma (\mathrm{T}) & \longrightarrow & \pi_ {1} (\mathrm{H}) & \longrightarrow & 0 \\ & & \Big \downarrow & & \Big \downarrow & & \Big \downarrow & & \\ 0 & \longrightarrow & \mathrm{N(G,T)} & \longrightarrow & \varGamma (\mathrm{T}) & \longrightarrow & \pi_ {1} (\mathrm{G}) & \longrightarrow & 0. \end{array}
\]

Remark. It can be shown (cf.~Exercise 2 of §5) that \(\pi_2(\mathrm{G / H})\) is zero. The exactness of the preceding sequence then gives an isomorphism from \(\pi_2(\mathrm{G / H})\) to \(\mathrm{N(G,T) / N(H,T)}\) .

COROLLARY 3. The homomorphism \(\pi_1(\mathrm{D}(\mathrm{G})) \to \pi_1(\mathrm{G})\) corresponding to the inclusion of \(\mathrm{D}(\mathrm{G})\) into \(\mathrm{G}\) induces an isomorphism from \(\pi_1(\mathrm{D}(\mathrm{G}))\) to the torsion subgroup of \(\pi_1(\mathrm{G})\) .

Indeed, \(\mathrm{T} \cap \mathrm{D}(\mathrm{G})\) is a maximal torus of \(\mathrm{D}(\mathrm{G})\) ( \(\S 2\) , no. 3, Prop. 1 \(c\) ); from the exact sequence

\[
0 \longrightarrow \Gamma (\mathrm{T} \cap \mathrm{D} (\mathrm{G})) \longrightarrow \Gamma (\mathrm{T}) \longrightarrow \Gamma (\mathrm{T} / (\mathrm{T} \cap \mathrm{D} (\mathrm{G}))) \longrightarrow 0
\]

and Proposition 11, we obtain an exact sequence

\[
0 \longrightarrow \pi_ {1} (\mathrm{D} (\mathrm{G})) \longrightarrow \pi_ {1} (\mathrm{G}) \longrightarrow \Gamma (\mathrm{T} / (\mathrm{T} \cap \mathrm{D} (\mathrm{G}))) \longrightarrow 0,
\]

hence the corollary, since \(\pi_{1}(\mathrm{D}(\mathrm{G}))\) is finite and \(\Gamma(\mathrm{T}/(\mathrm{T}\cap\mathrm{D}(\mathrm{G})))\) is free.

\subsection{SUBGROUPS OF MAXIMAL RANK}

Recall (Chap. VI, §1, no. 7) that a subset P of R = R(G, T) is said to be closed if \((P + P) \cap R \subset P\) , and symmetric if P = -P.

PROPOSITION 12. Let H be the set of connected closed subgroups of G containing T, ordered by inclusion. The map \(\mathrm{H} \mapsto \mathrm{R}(\mathrm{H}, \mathrm{T})\) is an increasing bijection from H to the set of symmetric closed subsets of \(\mathrm{R}(\mathrm{G}, \mathrm{T})\) , ordered by inclusion.

If \(H \in H\) , then \(\mathrm{L}(\mathrm{H})_{(\mathbf{C})}\) is the direct sum of \(t_{C}\) and the \(g^{\alpha}\) for \(\alpha \in \mathrm{R}(\mathrm{H}, \mathrm{T})\) ; since this is a reductive subalgebra in \(g_{C}\) , the subset \(\mathrm{R}(\mathrm{H}, \mathrm{T})\) of R satisfies the stated conditions (Chap. VIII, §3, no. 1, Lemma 2 and Prop. 2). Conversely, if P is a subset of R satisfying these conditions, then \(t_{C} \oplus \sum_{\alpha \in P} g^{\alpha}\) is a subalgebra of \(g_{C}\) (loc. cit.) which is rational over R (no. 3), and hence of the form \(h_{(C)}\) , where h is a subalgebra of g. Let \(\mathrm{H}(\mathrm{P})\) be the integral subgroup of G defined by h; it is closed (§2, no. 4, Remark 1). We verify immediately that the maps \(\mathrm{H} \mapsto \mathrm{R}(\mathrm{H}, \mathrm{T})\) and \(\mathrm{P} \mapsto \mathrm{H}(\mathrm{P})\) are increasing and inverses of each other.

COROLLARY 1. There are only finitely-many closed subgroups of G containing T.

Let H be such a subgroup; then \(H_{0} \in H\) , and H is finite. Moreover, H is a subgroup of \(\mathrm{N}_{\mathrm{G}}(\mathrm{H}_{0})\) containing \(H_{0}\) , and \(\mathrm{N}_{\mathrm{G}}(\mathrm{H}_{0})/\mathrm{H}_{0}\) is finite ( \(\S2\) , no. 4, Prop. 4 and Remark 2).

COROLLARY 2. Let H be a connected closed subgroup of G containing T, and let \(W_{\mathrm{G}}^{\mathrm{H}}(\mathrm{T})\) be the stabilizer in \(W_{\mathrm{G}}(\mathrm{T})\) of the subset \(\mathrm{R}(\mathrm{H},\mathrm{T})\) of R. The group \(N_{\mathrm{G}}(\mathrm{H})/\mathrm{H}\) is isomorphic to the quotient group \(W_{\mathrm{G}}^{\mathrm{H}}(\mathrm{T})/W_{\mathrm{H}}(\mathrm{T})\) .

Indeed, it follows from Prop. 7 of §2, no. 5, applied to \(\mathrm{N}_{\mathrm{G}}(\mathrm{H})\) , that \(\mathrm{N}_{\mathrm{G}}(\mathrm{H})/\mathrm{H}\) is isomorphic to \(\mathrm{W}_{\mathrm{N}(\mathrm{H})}(\mathrm{T})/\mathrm{W}_{\mathrm{H}}(\mathrm{T})\) , where \(\mathrm{W}_{\mathrm{N}(\mathrm{H})}(\mathrm{T})\) is the set of elements of \(\mathrm{W}_{\mathrm{G}}(\mathrm{T})\) whose representatives in \(\mathrm{N}_{\mathrm{G}}(\mathrm{T})\) normalize H. Let \(n \in \mathrm{N}_{\mathrm{G}}(\mathrm{T})\) , and let w be its class in \(\mathrm{W}_{\mathrm{G}}(\mathrm{T})\) . By Chap. III, §9, no. 4, Prop. 11, n normalizes H if and only if \((\operatorname{Ad} n)(\operatorname{L}(\operatorname{H})) = \operatorname{L}(\operatorname{H})\) ; in view of Prop. 5 of no. 3, this also means that the subset \(\operatorname{R}(\operatorname{H}, \operatorname{T})\) of R is stable under w, hence the corollary.

Remark 1. The group \(W_{\mathrm{G}}^{\mathrm{H}}(\mathrm{T})\) is also the stabilizer in \(W_{\mathrm{G}}(\mathrm{T})\) of the subgroup \(\mathrm{C}(\mathrm{H})\) of T: this follows from Prop. 8 of no. 4.

PROPOSITION 13. Let H be a connected closed subgroup of G of maximal rank, and C its centre. Then C contains the centre of G, and H is the identity component of the centralizer of C.

Let S be a maximal torus of H. Since the centre of G is contained in S, it is contained in C. Put \(L = Z(C)_{0}\) ; this is a connected closed subgroup of G containing H, hence is of maximal rank, and its centre is equal to C. Denote by \(R_{H}\) and \(R_{L}\) the root systems of H and L, respectively, relative to S; then \(R_{H} \subset R_{L} \subset R(G, S)\) . Since \(C(H) = C(L)\) , Prop. 8 (no. 4) implies the equality \(Q(R_{H}) = Q(R_{L})\) ; but \(Q(R_{H}) \cap R_{L} = R_{H}\) (Chap. VI, §1, no. 7, Prop. 23), so \(R_{H} = R_{L}\) and H = L (Prop. 12).

Remark 2. Say that a subgroup \(\mathrm{C}\) of \(\mathrm{G}\) is radical if there exists a maximal torus \(\mathrm{S}\) of \(\mathrm{G}\) and a subset \(\mathrm{P}\) of \(\mathrm{R}(\mathrm{G},\mathrm{S})\) such that \(\mathrm{C} = \bigcap_{\alpha \in \mathrm{P}}\mathrm{Ker}\alpha\) . It follows

from Prop. 13 and Lemma 2 of no. 5 that the map \(H \mapsto C(H)\) induces a bijection from the set of connected closed subgroups of maximal rank to the set of radical subgroups of G. The inverse bijection is the map \(C \mapsto Z(C)_{0}\) .

COROLLARY. The set of \(g \in \mathbf{G}\) such that \(\mathrm{T} \cap g\mathrm{T}g^{-1} \neq \mathrm{C}(\mathrm{G})\) is the union of a finite number of closed analytic submanifolds of \(\mathbf{G}\) distinct from \(\mathbf{G}\) .

Indeed, put \(A_{g} = T \cap gTg^{-1}\) ; we have \(T \subset Z(A_{g})\) and \(gTg^{-1} \subset Z(A_{g})\) . Hence, there exists \(x \in Z(A_{g})\) such that \(xTx^{-1} = gTg^{-1}\) (§2, no. 2, Th. 2), which implies that \(g \in Z(A_{g}).N_{G}(T)\) . Denote by A the finite (Cor. 1) set of closed subgroups of G containing T and distinct from G, and put \(X = \bigcup_{H \in A} H.N_{G}(T)\) ; this is a finite union of closed submanifolds of G, distinct from G. If \(A_{g} \neq C(G)\) , then \(Z(A_{g}) \in A\) , and g belongs to X. Conversely, if \(g \in H.N_{G}(T)\) , with \(H \in A\) , then \(A_{g}\) contains \(C(H)\) , so \(A_{g} \neq C(G)\) (Prop. 13).

PROPOSITION 14. Let X be a subset of T, and let \(R_{X}\) be the set of roots \(\alpha \in \mathrm{R}(\mathrm{G}, \mathrm{T})\) such that \(\alpha(\mathrm{X}) = \{1\}\) . The group \(\mathrm{Z}_{\mathrm{G}}(\mathrm{X}) / \mathrm{Z}_{\mathrm{G}}(\mathrm{X})_{0}\) is isomorphic to the quotient of the subgroup of \(\mathrm{W}_{\mathrm{G}}(\mathrm{T})\) fixing X by the subgroup generated by the reflections \(s_{\alpha}\) for \(\alpha \in R_{X}\) .

Put \(H = Z_{\mathrm{G}}(\mathrm{X})\) ; since \(\mathrm{L}(\mathrm{H})_{(\mathbf{C})}\) is the set of points of \(g_{C}\) fixed by \(\mathrm{Ad}(\mathrm{X})\) , it is the sum of \(t_{C}\) and the \(g^{\alpha}\) for which \(\alpha(\mathrm{X}) = \{1\}\) . Consequently, \(\mathrm{R}(\mathrm{H}_{0}, \mathrm{T}) = \mathrm{R}_{\mathrm{X}}\) , so \(\mathrm{W}_{\mathrm{H}_{0}}(\mathrm{T})\) is generated by the reflections \(s_{\alpha}\) for \(\alpha \in R_{X}\) . It now suffices to apply Prop. 7 of §2, no. 5.

We shall see below ( \(\S5\) , no. 3, Th. 1) that if G is simply-connected and X reduces to a point, the centralizer \(Z(X)\) is connected.

\subsection{ROOT DIAGRAMS}

DEFINITION 2. A root diagram (or simply a diagram, if there is no risk of confusion) is a triple \(D = (M, M_{0}, R)\) where:

\((RD_{0})\) M is a free Z-module of finite type and the submodule \(M_{0}\) is a direct factor of M;

\((\mathrm{RD}_{\mathrm{I}})\) R is a finite subset of M; \(\mathrm{R} \cup \mathrm{M}_0\) generates the Q-vector space \(\mathbf{Q} \otimes \mathrm{M}\) ;

\((\mathrm{RD}_{\mathrm{II}})\) for all \(\alpha \in \mathrm{R}\) , there exists an element \(\alpha^{\vee}\) of \(\mathrm{M}^{*} = \mathrm{Hom}_{\mathbf{Z}}(\mathrm{M}, \mathbf{Z})\) such that \(\alpha^{\vee}(\mathrm{M}_{0}) = 0\) , \(\alpha^{\vee}(\alpha) = 2\) and the endomorphism \(x \mapsto x - \alpha^{\vee}(x)\alpha\) of M leaves R stable.

By Chap. VI, §1, no. 1, for all \(\alpha \in \mathbb{R}\) the element \(\alpha^{\vee}\) of \(\mathrm{M}^*\) is uniquely determined by \(\alpha\) ; we denote by \(s_{\alpha}\) the endomorphism \(x \mapsto x - \alpha^{\vee}(x)\alpha\) of M. Moreover (loc. cit.), the \(\mathbf{Q}\) -vector space \(\mathbf{Q} \otimes \mathrm{M}\) is the direct sum of \(\mathbf{Q} \otimes \mathrm{M}_0\) and the vector subspace V(R) generated by R, and R is a root system in V(R) (loc. cit., Def. 1).

The elements of R are called the roots of the root diagram D, and the elements \(\alpha^{\vee}\) of \(M^{*}\) the inverse roots. The group generated by the automorphisms \(s_{\alpha}\) of M is called the Weyl group of D and is denoted by \(\mathrm{W(D)}\) ; the elements of \(\mathrm{W(D)}\) induce the identity on \(M_{0}\) , and induce on \(\mathrm{V(R)}\) the transformations of the Weyl group of the root system R.

Examples. 1) For every free \(\mathbf{Z}\) -module of finite type \(M\) , the triple \((M, M, \varnothing)\) is a root diagram.

\begin{enumerate}
\def\labelenumi{\arabic{enumi})}
\setcounter{enumi}{1}
\tightlist
\item
  If \(D = (M, M_{0}, R)\) is a root diagram, let \(M_{0}^{*}\) be the orthogonal complement of \(V(R)\) in \(M^{*}\) , and let \(R^{\vee}\) be the set of inverse roots of D. Then \(D^{\vee} = (M^{*}, M_{0}^{*}, R^{\vee})\) is a root diagram, called the inverse of D. For all \(\alpha \in R\) , the symmetry \(s_{\alpha^{\vee}}\) of \(M^{*}\) is the contragredient automorphism of the symmetry \(s_{\alpha}\) of M; the map \(w \mapsto {}^{t} w^{-1}\) is an isomorphism from \(W(D)\) to \(W(D^{\vee})\) . Moreover, \(V(R^{\vee})\) can be naturally identified with the dual of the Q-vector space \(V(R)\) , \(R^{\vee}\) then being identified with the inverse root system of R.
\end{enumerate}

If the dual of \(\mathbf{M}^*\) is identified with M, the inverse diagram of \(\mathrm{D}^{\vee}\) is identified with D.

\begin{enumerate}
\def\labelenumi{\arabic{enumi})}
\setcounter{enumi}{2}
\item
  Let \((\mathfrak{g},\mathfrak{h})\) be a split reductive \(\mathbf{Q}\) -Lie algebra, and \(\mathrm{M} \subset \mathfrak{h}\) a permissible lattice (Chap. VIII, §12, no. 6, Def. 1). Let \(\mathrm{M}_0\) be the subgroup of M orthogonal to the roots of \((\mathfrak{g},\mathfrak{h})\) and \(\mathrm{R}^\vee\) the set of the \(H_{\alpha}\) , \(\alpha \in \mathrm{R}(\mathfrak{g},\mathfrak{h})\) . Then \((\mathrm{M},\mathrm{M}_0,\mathrm{R}^\vee)\) is a root diagram, and \((\mathrm{M}^*,\mathrm{M}_0^*,\mathrm{R}(\mathfrak{g},\mathfrak{h}))\) is the inverse diagram.
\item
  Let V be a vector space over Q and R a root system in V; denote by P(R) the group of weights of R and by Q(R) the group of radical weights of R (Chap. VI, §1, no. 9). Then (Q(R), 0, R) and (P(R), 0, R) are root diagrams. A root diagram (M, M₀, S) is isomorphic to a diagram of the form (Q(R), 0, R) (resp. (P(R), 0, R)) if and only if M is generated by S (resp. M* is generated by S \(^{v}\) ).
\end{enumerate}

For every subgroup X of P(R) containing Q(R), (X, 0, R) is a root diagram and, up to isomorphism, every diagram (M, M₀, S) such that M₀ = 0, in other words such that S generates a subgroup of M of finite index, arises in this way.

The root diagram \((\mathrm{M},\mathrm{M}_0,\mathrm{R})\) is said to be reduced if the root system \(\mathbf{R}\) is reduced (in other words (Chap. VI, §1, no. 4) if the relations \(\alpha ,\beta \in \mathbb{R}\) , \(\lambda \in \mathbf{Z}\) , \(\beta = \lambda \alpha\) imply that \(\lambda = 1\) or \(\lambda = -1\) ). The diagrams in Examples 1) and 3) are reduced.

\subsection{COMPACT LIE GROUPS AND ROOT SYSTEMS}

With the terminology introduced in the preceding number, an important part of the results of numbers 4 and 5 can be summarized in the following theorem:

THEOREM 2. a) (X(T), X(T/(T ∩ D(G)), R(G, T))) is a reduced root diagram; its Weyl group consists of the X(w) for w ∈ W; the group X(C(G)) is isomorphic to the quotient of X(T) by the subgroup generated by R(G, T).

\begin{enumerate}
\def\labelenumi{\alph{enumi})}
\setcounter{enumi}{1}
\item
  \((\Gamma(\mathrm{T}),\Gamma(\mathrm{C}(\mathrm{G})_{0}),\mathrm{R}^{\vee}(\mathrm{G},\mathrm{T}))\) is a reduced root diagram; its Weyl group consists of the \(\Gamma(w)\) , for \(w\in W\) ; the group \(\pi_{1}(\mathrm{G})\) is isomorphic to the quotient of \(\Gamma(\mathrm{T})\) by the subgroup generated by \(\mathrm{R}^{\vee}(\mathrm{G},\mathrm{T})\) .
\item
  If each of the Z-modules X(T) and \(\Gamma(T)\) is identified with the dual of the other (no. 2, Prop. 3), each of the preceding root diagrams is identified with the inverse of the other.
\end{enumerate}

Denote by \(\mathrm{D}^{*}(\mathrm{G},\mathrm{T})\) the diagram \((\mathrm{X}(\mathrm{T}),\mathrm{X}(\mathrm{T}/(\mathrm{T}\cap\mathrm{D}(\mathrm{G})),\mathrm{R}(\mathrm{G},\mathrm{T})))\) and by \(\mathrm{D}_{*}(\mathrm{G},\mathrm{T})\) the diagram \((\varGamma(\mathrm{T}),\varGamma(\mathrm{C}(\mathrm{G})_{0}),\mathrm{R}^{\vee}(\mathrm{G},\mathrm{T}))\) ; these are called the contravariant diagram and the covariant diagram of G (relative to T), respectively.

Examples. 1) If G is semi-simple of rank 1, then \(\mathrm{D}^{*}(\mathrm{G},\mathrm{T})\) and \(\mathrm{D}_{*}(\mathrm{G},\mathrm{T})\) are necessarily isomorphic to one of the two diagrams \(\Delta_{2}=(\mathbf{Z},0,\{2,-2\})\) , \(\Delta_{1}=(\mathbf{Z},0,\{1,-1\})\) . If G is isomorphic to \(\mathbf{SU}(2,\mathbf{C})\) , \(\mathrm{D}_{*}(\mathrm{G},\mathrm{T})\) is isomorphic to \(\Delta_{1}\) (since G is simply-connected) so \(\mathrm{D}^{*}(\mathrm{G},\mathrm{T})\) is isomorphic to \(\Delta_{2}\) . If G is isomorphic to \(\mathbf{SO}(3,\mathbf{R})\) , \(\mathrm{D}^{*}(\mathrm{G},\mathrm{T})\) is isomorphic to \(\Delta_{1}\) (since \(\mathrm{C}(\mathrm{G})=\{1\})\) , so \(\mathrm{D}_{*}(\mathrm{G},\mathrm{T})\) is isomorphic to \(\Delta_{2}\) .

\begin{enumerate}
\def\labelenumi{\arabic{enumi})}
\setcounter{enumi}{1}
\item
  If G and \(G'\) are two connected compact Lie groups, with maximal tori T and \(T'\) , respectively, and if \(\mathrm{D}^{*}(\mathrm{G},\mathrm{T})=(\mathrm{M},\mathrm{M}_{0},\mathrm{R})\) and \(\mathrm{D}^{*}(\mathrm{G}',\mathrm{T}')=(\mathrm{M}',\mathrm{M}_{0}',\mathrm{R}')\) , then \(\mathrm{D}^{*}(\mathrm{G}\times\mathrm{G}',\mathrm{T}\times\mathrm{T}')\) can be identified with \((\mathrm{M}\oplus\mathrm{M}',\mathrm{M}_{0}\oplus\mathrm{M}_{0}',\mathrm{R}\cup\mathrm{R}')\) . Similarly for the covariant diagrams.
\item
  Let N be a closed subgroup of T, central in G, and let \((M, M_{0}, R)\) be the contravariant diagram of G relative to T. Then the contravariant diagram of G/N relative to T/N can be identified with \((M', M_{0}', R)\) , where \(M'\) is the subgroup of M consisting of the \(\lambda\) such that \(\lambda(N) = \{1\}\) and \(M_{0}' = M' \cap M_{0}\) .
\item
  Similarly, let N be a finite abelian group, and \(\varphi : \pi_{1}(G) \to N\) a surjective homomorphism. Let \(G'\) be the covering of G associated to this homomorphism; this is a connected compact Lie group, of which N is a central subgroup (General Topology, Chap. XI, in preparation), and G can be naturally identified with \(G'/N\) . Let \(T'\) be the maximal torus of \(G'\) that is the inverse image of T. If \((P, P_{0}, S)\) is the covariant diagram of G relative to T, the covariant diagram of \(G'\) relative to \(T'\) can be identified with \((P', P_{0}', S)\) , where \(P'\) is the kernel of the composite homomorphism \(\varphi \circ f(G, T) : P \to N\) (cf.~no. 6, Prop. 11), and \(P_{0}' = P_{0} \cap P'\) .
\end{enumerate}

Remarks. 1) Let \(\mathfrak{c}\) be the centre of \(\mathfrak{g}_{\mathbf{C}}\) ; then \(\mathfrak{c} = \mathrm{L}(\mathrm{C}(\mathrm{G}))_{(\mathbf{C})}\) . We have the following relations between the diagrams of G relative to T and the direct and inverse root systems of the split reductive algebra \((\mathfrak{g}_{\mathbf{C}}, t_{\mathbf{C}})\) :

\begin{enumerate}
\def\labelenumi{\alph{enumi})}
\tightlist
\item
  The canonical isomorphism from \(\mathbf{C} \otimes \Gamma(\mathrm{T})\) to \(\mathfrak{t}_{\mathbf{C}}\) induces a bijection from \(\mathbf{C} \otimes \Gamma(\mathrm{C}(\mathrm{G})_0)\) to \(\mathfrak{c}\) and a bijection from \(1 \otimes \mathrm{R}^{\vee}(\mathrm{G}, \mathrm{T})\) to \(2\pi i.\mathrm{R}^{\vee}(\mathfrak{g}_{\mathbf{C}}, \mathfrak{t}_{\mathbf{C}})\) .
\end{enumerate}

from \(\mathbf{C}\otimes\varGamma(\mathrm{C}(\mathrm{G})_{0})\) to c and a bijection from \(1\otimes\mathrm{R}^{\vee}(\mathrm{G},\mathrm{T})\) to \(2\pi i.\mathrm{R}^{\vee}(\mathfrak{g}_{\mathbf{C}},\mathfrak{t}_{\mathbf{C}})\) . b) The canonical isomorphism from \(\mathbf{C}\otimes\mathrm{X}(\mathrm{T})\) to the dual \(t_{C}^{*}\) of \(t_{C}\) induces a bijection from \(\mathbf{C}\otimes\mathrm{X}(\mathrm{T}/(\mathrm{T}\cap\mathrm{D}(\mathrm{G})))\) to the orthogonal complement of \(t_{C}\cap\mathcal{D}(\mathfrak{g})_{C}\) , and a bijection from \(1\otimes\mathrm{R}(\mathrm{G},\mathrm{T})\) to \(\mathrm{R}(\mathfrak{g}_{\mathbf{C}},\mathfrak{t}_{\mathbf{C}})\) .

\begin{enumerate}
\def\labelenumi{\arabic{enumi})}
\setcounter{enumi}{1}
\tightlist
\item
  Assume that the group G is semi-simple; denote by R (resp. \(R^{\vee}\) ) the root system R(G, T) (resp. \(R^{\vee}(G, T)\) ), so that we have inclusions
\end{enumerate}

\[
\mathrm{Q} (\mathrm{R}) \subset \mathrm{X} (\mathrm{T}) \subset \mathrm{P} (\mathrm{R}) \quad \mathrm{Q} (\mathrm{R} ^ {\vee}) \subset \Gamma (\mathrm{T}) \subset \mathrm{P} (\mathrm{R} ^ {\vee}).
\]

The finite abelian groups \(\mathrm{P}(\mathrm{R})/\mathrm{Q}(\mathrm{R})\) and \(\mathrm{P}(\mathrm{R}^{\vee})/\mathrm{Q}(\mathrm{R}^{\vee})\) are in duality (Chap. VI, §1, no. 9); if \(M^{*}\) denotes the dual group of the finite abelian group M, we deduce from the preceding canonical isomorphisms

\[
\begin{array}{l l} \Gamma (\mathrm{T}) / \mathrm{Q} (\mathrm{R} ^ {\vee}) \to \pi_ {1} (\mathrm{G}) & \mathrm{P} (\mathrm{R} ^ {\vee}) / \Gamma (\mathrm{T}) \to \mathrm{C} (\mathrm{G}) \\ \mathrm{P} (\mathrm{R}) / \mathrm{X} (\mathrm{T}) \to (\pi_ {1} (\mathrm{G})) ^ {\wedge} & \mathrm{X} (\mathrm{T}) / \mathrm{Q} (\mathrm{R}) \to (\mathrm{C} (\mathrm{G})) ^ {\wedge}. \end{array}
\]

In particular, the product of the orders of \(\pi_{1}(\mathrm{G})\) and \(\mathrm{C}(\mathrm{G})\) is equal to the connection index f of \(\mathrm{R}(\mathrm{G},\mathrm{T})\) (loc. cit.).

Now let \(G'\) be another connected compact Lie group, \(T'\) a maximal torus of \(G'\) . Let \(f: G \to G'\) be an isomorphism of Lie groups such that \(f(T) = T'\) ; denote by \(f_{T}\) the isomorphism from T to \(T'\) that it defines. Then \(\mathrm{X}(f_{\mathrm{T}})\) is an isomorphism from \(\mathrm{D}^{*}(\mathrm{G}', \mathrm{T}')\) to \(\mathrm{D}^{*}(\mathrm{G}, \mathrm{T})\) , denoted by \(\mathrm{D}^{*}(f)\) , and \(\Gamma(f_{\mathrm{T}})\) is an isomorphism from \(\mathrm{D}_{*}(\mathrm{G}, \mathrm{T})\) to \(\mathrm{D}_{*}(\mathrm{G}', \mathrm{T}')\) , denoted by \(\mathrm{D}_{*}(f)\) . If \(t \in T\) , and if we put \(g = f \circ \operatorname{Int} t = (\operatorname{Int} f(t)) \circ f\) , then \(\mathrm{D}^{*}(g) = \mathrm{D}^{*}(f)\) , \(\mathrm{D}_{*}(g) = \mathrm{D}_{*}(f)\) .

PROPOSITION 15. Let \(\varphi\) be an isomorphism from \(\mathrm{D}^*(\mathrm{G}',\mathrm{T}')\) to \(\mathrm{D}^*(\mathrm{G},\mathrm{T})\) (resp. from \(\mathrm{D}_*(\mathrm{G},\mathrm{T})\) to \(\mathrm{D}_*(\mathrm{G}',\mathrm{T}')\) ). There exists an isomorphism \(f:\mathrm{G}\to \mathrm{G}'\) such that \(f(\mathrm{T}) = \mathrm{T}'\) and \(\varphi = \mathrm{D}^*(f)\) (resp. \(\varphi = \mathrm{D}_*(f)\) ); if \(f_{1}\) and \(f_{2}\) are two such isomorphisms, there exists an element \(t\) of \(\mathrm{T}\) such that \(f_{2} = f_{1}\circ \operatorname {Int}t\) .

The second assertion follows immediately from Prop. 9 (no. 4); we prove the first for the covariant diagrams, for example. Denote by \(\mathfrak{g}'\) (resp. \(\mathfrak{t}'\) ) the Lie algebra of \(G'\) (resp. \(T'\) ), and by \(\mathfrak{g}_{\mathbf{C}}'\) (resp. \(\mathfrak{t}_{\mathbf{C}}'\) ) its complexified Lie algebra. By Chap. VIII, §4, no. 4, Th. 2 (i), there exists an isomorphism \(\psi : \mathfrak{g}_{\mathbf{C}} \to \mathfrak{g}_{\mathbf{C}}'\) that maps \(\mathfrak{t}_{\mathbf{C}}\) to \(\mathfrak{t}_{\mathbf{C}}'\) and induces on \(\Gamma(T) \subset \mathfrak{t}_{\mathbf{C}}\) the given isomorphism \(\varphi : \Gamma(T) \to \Gamma(T')\) . Then \(\mathfrak{g}\) and \(\psi^{-1}(\mathfrak{g}')\) are two compact forms of \(\mathfrak{g}_{\mathbf{C}}\) that have the same intersection \(\mathfrak{t}\) with \(\mathfrak{t}_{\mathbf{C}}\) ; by §3, no. 2, Prop. 3, there exists an inner automorphism \(\theta\) of \(\mathfrak{g}_{\mathbf{C}}\) inducing the identity on \(\mathfrak{t}_{\mathbf{C}}\) and such that \(\theta(\mathfrak{g}) = \psi^{-1}(\mathfrak{g}')\) . By replacing \(\psi\) by \(\psi \circ \theta\) , we can assume that \(\psi\) maps \(\mathfrak{g}\) to \(\mathfrak{g}'\) . Further, by Prop. 4 of no. 2, there exists a unique morphism \(f_{\mathrm{T}} : T \to T'\) such that \(\Gamma(f_{\mathrm{T}}) = \varphi\) . Then the restriction of \(\psi\) to \(\mathfrak{t}\) is \(L(f_{\mathrm{T}})\) , and by §2, no. 6, Prop. 8, there exists a unique morphism \(f : G \to G'\) that induces \(f_{\mathrm{T}}\) on \(T\) and \(\psi\) on \(g_{C}\) . Applying the preceding to \(\varphi^{-1}\) and \(\psi^{-1}\) we obtain an inverse morphism to f, which is therefore an isomorphism. Then \(\mathrm{D}_{*}(f)=\Gamma(f_{\mathrm{T}})=\varphi\) , hence the proposition.

Note that, if T and \(T'\) are two maximal tori of G, the diagrams \(\mathrm{D}^{*}(\mathrm{G},\mathrm{T})\) and \(\mathrm{D}^{*}(\mathrm{G},\mathrm{T}')\) are isomorphic (if \(g \in G\) is such that \(gTg^{-1} = T'\) , then \(\operatorname{Int} g\) is an isomorphism from G to G that maps T to \(T'\) ). Denote by \(\mathrm{D}^{*}(\mathrm{G})\) the isomorphism class of \(\mathrm{D}^{*}(\mathrm{G},\mathrm{T})\) (cf.~Theory of Sets, Chap. II, §6, no. 2); this is a root diagram that depends only on G and is called the contravariant diagram of G. The covariant diagram \(\mathrm{D}_{*}(\mathrm{G})\) of G is defined similarly, and we obtain:

COROLLARY. Two connected compact Lie groups G and \(G'\) are isomorphic if and only if the diagrams \(\mathrm{D}^{*}(\mathrm{G})\) and \(\mathrm{D}^{*}(\mathrm{G}')\) (resp. \(\mathrm{D}_{*}(\mathrm{G})\) and \(\mathrm{D}_{*}(\mathrm{G}')\) ) are equal.

PROPOSITION 16. For every reduced root diagram D, there exists a connected compact Lie group G such that \(\mathrm{D}^{*}(\mathrm{G})\) (resp. \(\mathrm{D}_{*}(\mathrm{G})\) ) is isomorphic to D.

\begin{enumerate}
\def\labelenumi{\alph{enumi})}
\item
  By replacing D, if necessary, by its inverse diagram, we are reduced to constructing G such that \(\mathrm{D}^{*}(\mathrm{G})\) is isomorphic to D. Put \(\mathrm{D} = (\mathrm{M}, \mathrm{M}_{0}, \mathrm{R})\) ; then \(Q \otimes M\) is the direct sum of \(Q \otimes M_{0}\) and the vector subspace \(\mathrm{V}(\mathrm{R})\) generated by R. Moreover, since the inverse roots take integer values on M, the projection of M on \(\mathrm{V}(\mathrm{R})\) parallel to \(Q \otimes M_{0}\) is contained in the group of weights \(\mathrm{P}(\mathrm{R})\) of R, so that M is a subgroup of \(\mathrm{M}_{0} \oplus \mathrm{P}(\mathrm{R})\) of finite index. Denote by \(D'\) the diagram \((\mathrm{M}_{0} \oplus \mathrm{P}(\mathrm{R}), \mathrm{M}_{0}, \mathrm{R})\) .
\item
  Let \(\mathfrak{a}\) be a complex semi-simple Lie algebra whose canonical root system is isomorphic to \(R \subset \mathbf{C} \otimes V(R)\) (Chap. VIII, §4, no. 3), and let \(\mathfrak{g}_1\) be a compact real form of \(\mathfrak{a}\) ( \(\S 3\) , no. 2, Th. 1). Let \(G_1\) be a simply-connected real Lie group whose Lie algebra is isomorphic to \(\mathfrak{g}_1\) ; then \(G_1\) is compact ( \(\S 1\) , no. 4, Th. 1). Let \(T_1\) be a maximal torus of \(G_1\) . By Th. 1, the diagram \(D^*(G_1, T_1)\) is isomorphic to \((P(R), 0, R)\) .
\item
  Let \(T_{0}\) be a torus of dimension equal to the rank of \(M_{0}\) ; then \(\mathrm{D}^{*}(\mathrm{T}_{0},\mathrm{T}_{0})\) is isomorphic to \((\mathrm{M}_{0},\mathrm{M}_{0},\varnothing)\) , so \(\mathrm{D}^{*}(\mathrm{G}_{1}\times\mathrm{T}_{0},\mathrm{T}_{1}\times\mathrm{T}_{0})\) is isomorphic to \(D'\) (Example 2).
\item
  Finally, let N be the finite subgroup of \(T_{1} \times T_{0}\) orthogonal to M. Put \(\mathrm{G} = (\mathrm{G}_{1} \times \mathrm{T}_{0})/\mathrm{N}\) , \(\mathrm{T} = (\mathrm{T}_{1} \times \mathrm{T}_{0})/\mathrm{N}\) . Then G is a connected compact Lie group, T a maximal torus of G, and D(G, T) is isomorphic to D (Example 3).
\end{enumerate}

Scholium. The classification of connected compact Lie groups up to isomorphism is thus reduced to that of reduced root diagrams. The connected compact semi-simple Lie groups correspond to the reduced root diagrams \((M, M_{0}, R)\) such that \(M_{0} = 0\) ; giving such a diagram is equivalent to giving a reduced root system R in a vector space V over Q and a subgroup M of V such that \(Q(R) \subset M \subset P(R)\) .

Remark 3. Let \(\mathrm{T}'\) be another maximal torus of G, B (resp. \(\mathrm{B}'\) ) a basis of the root system \(\mathrm{R}(\mathrm{G},\mathrm{T})\) (resp. \(\mathrm{R}(\mathrm{G}',\mathrm{T}')\) ) (Chap. VI, §1, no. 5, Def. 2). There exist elements \(g\in \mathrm{G}\) such that Int \(g\) maps T onto \(\mathrm{T}'\) and B onto \(\mathrm{B}'\) , and these elements form a unique coset modulo Int(T) (since T and \(\mathrm{T}'\) are conjugate, we can assume that \(\mathrm{T} = \mathrm{T}'\) , and it suffices to apply Chap. VI, §1, no. 5, Remark 4 and Prop. 9 of no. 4). It follows that the isomorphism from T to \(\mathrm{T}'\) induced by Int \(g\) is independent of the choice of \(g\) ; consequently the same is true of \(\mathrm{D}_{*}(\mathrm{Int}g)\) and \(\mathrm{D}^{*}(\mathrm{Int}g)\) . Paraphrasing Chap. VIII, §5, no. 3, Remark 2, mutatis mutandis, we can now define the canonical maximal torus of G, the canonical covariant and contravariant root diagrams of G, \ldots.

\subsection{AUTOMORPHISMS OF A CONNECTED COMPACT LIE GROUP}

Denote by \(\mathrm{Aut}(\mathrm{G})\) the Lie group of automorphisms of G (Chap. III, §10, no. 2), and by \(\mathrm{Aut}(\mathrm{G},\mathrm{T})\) the closed subgroup of \(\mathrm{Aut}(\mathrm{G})\) consisting of the elements \(u\) such that \(u(\mathrm{T}) = \mathrm{T}\) . We have seen (§1, no. 4, Cor. 5 of Prop. 4) that the identity component of \(\mathrm{Aut}(\mathrm{G})\) is the subgroup \(\mathrm{Int}(\mathrm{G})\) of inner automorphisms; denote by \(\mathrm{Int}_{\mathrm{G}}(\mathrm{H})\) the image in \(\mathrm{Int}(\mathrm{G})\) of a subgroup H of G.

Let D be the covariant diagram of G relative to T; denote by \(\operatorname{Aut}(D)\) the group of its automorphisms, and by \(\mathrm{W}(D)\) its Weyl group. The map \(u \mapsto \mathrm{D}_{*}(u)\) is a homomorphism from \(\operatorname{Aut}(G, T)\) to \(\operatorname{Aut}(D)\) . Prop. 15 of no. 9 immediately gives:

PROPOSITION 17. The homomorphism \(\operatorname{Aut}(\mathrm{G},\mathrm{T})\to \operatorname{Aut}(\mathrm{D})\) is surjective, with kernel \(\operatorname{Int}_{\mathrm{G}}(\mathrm{T})\) .

Note that \(\operatorname{Aut}(G,T) \cap \operatorname{Int}(G) = \operatorname{Int}_{G}(N_{G}(T))\) and that the image of \(\operatorname{Int}_{G}(N_{G}(T))\) in \(\operatorname{Aut}(D)\) is \(W(D)\) (no. 5, Prop. 10). Thus, Proposition 17 gives an isomorphism

\[
\operatorname{Aut} (\mathrm{G}, \mathrm{T}) / (\operatorname{Aut} (\mathrm{G}, \mathrm{T}) \cap \operatorname{Int} (\mathrm{G})) \rightarrow \operatorname{Aut} (\mathrm{D}) / \mathrm{W} (\mathrm{D}).
\]

Further, \(\operatorname{Aut}(\mathrm{G}) = \operatorname{Int}(\mathrm{G}).\operatorname{Aut}(\mathrm{G}, \mathrm{T})\) . Indeed, if u belongs to \(\operatorname{Aut}(\mathrm{G})\) , \(u(\mathrm{T})\) is a maximal torus of T, hence is conjugate to T, and there exists an inner automorphism v of G such that \(u(\mathrm{T}) = v(\mathrm{T})\) , in other words \(v^{-1}u \in \operatorname{Aut}(\mathrm{G}, \mathrm{T})\) . It follows that \(\operatorname{Aut}(\mathrm{G})/\operatorname{Int}(\mathrm{G})\) can be identified with \(\operatorname{Aut}(\mathrm{G}, \mathrm{T})/(\operatorname{Aut}(\mathrm{G}, \mathrm{T}) \cap \operatorname{Int}(\mathrm{G}))\) , so in view of the preceding we have an exact sequence

\[
1 \to \operatorname{Int} (\mathrm{G}) \to \operatorname{Aut} (\mathrm{G}) \to \operatorname{Aut} (\mathrm{D}) / \mathrm{W} (\mathrm{D}) \to 1.\tag{16}
\]

Consequently:

PROPOSITION 18. The group \(\operatorname{Aut}(\mathrm{G}) / \operatorname{Int}(\mathrm{G})\) is isomorphic to \(\operatorname{Aut}(\mathrm{D}) / \operatorname{W}(\mathrm{D})\) .

In particular, assume that G is semi-simple; the group \(\operatorname{Aut}(D)\) can then be identified with the subgroup of \(\mathrm{A}(\mathrm{R}(\mathrm{G},\mathrm{T}))\) (Chap. VI, §1, no. 1) consisting of the elements u such that \(u(\mathrm{X}(\mathrm{T})) \subset \mathrm{X}(\mathrm{T})\) , and the subgroup \(\mathrm{W}(\mathrm{D})\) can be identified with \(\mathrm{W}(\mathrm{R}(\mathrm{G},\mathrm{T}))\) .

COROLLARY. If G is simply-connected, or if C(G) reduces to the identity element, the group \(\operatorname{Aut}(\mathrm{G})/\operatorname{Int}(\mathrm{G})\) is isomorphic to the group of automorphisms of the Dynkin graph of R(G,T).

This follows from the preceding and Chap. VI, §4, no. 2, Cor. of Prop. 1.

We are now going to show that the extension (16) admits sections.

For all \(\alpha\in\mathrm{R}(\mathrm{G},\mathrm{T})\) , denote by \(\mathrm{V}(\alpha)\) the 2-dimensional vector subspace of g such that \(\mathrm{V}(\alpha)_{(\mathbf{C})}=\mathfrak{g}^{\alpha}+\mathfrak{g}^{-\alpha}\) ; denote by K the quadratic form associated to the Killing form of g.

DEFINITION 3. A framing of (G,T) is a pair \((\mathrm{B},(U_{\alpha})_{\alpha \in \mathrm{B}})\) , where B is a basis of \(\mathrm{R}(\mathrm{G},\mathrm{T})\) (Chap. VI, §1, no. 5, Def. 2) and where, for all \(\alpha \in \mathrm{B}\) , \(U_{\alpha}\) is an element of \(\mathrm{V}(\alpha)\) such that \(\mathrm{K}(U_{\alpha}) = -1\) .

A framing of G is a maximal torus T of G together with a framing of \((G, T)\) .

Lemma 3. Let \(B_{0}\) be a basis of \(\mathrm{R}(\mathrm{G},\mathrm{T})\) . The group \(\operatorname{Int}_{\mathrm{G}}(\mathrm{T})\) operates simply-transitively on the set of framings of \((\mathrm{G},\mathrm{T})\) of the form \((\mathrm{B}_{0},(U_{\alpha})_{\alpha\in\mathrm{B}_{0}})\) .

For all \(\alpha\in B_{0}\) , denote by \(\mathrm{K}(\alpha)\) the restriction of the quadratic form K to \(\mathrm{V}(\alpha)\) ; the operation of T on \(\mathrm{V}(\alpha)\) defines a morphism \(\iota_{\alpha}:T\to\mathbf{SO}(\mathrm{K}(\alpha))\) . We have seen in no. 4 that \(\mathbf{SO}(\mathrm{K}(\alpha))\) can be identified with U in such a way that \(\iota_{\alpha}\) is identified with the root \(\alpha\) . Since \(B_{0}\) is a basis of R, it is a basis of the Z-module \(\mathrm{Q}(\mathrm{R})\) generated by the roots, hence a basis of the submodule \(\mathrm{X}(\mathrm{T}/\mathrm{C}(\mathrm{G}))\) of \(\mathrm{X}(\mathrm{T})\) . It follows that the product of the morphisms \(\iota_{\alpha}\) induces an isomorphism from \(\mathrm{T}/\mathrm{C}(\mathrm{G})\) to the product of the groups \(\mathbf{SO}(\mathrm{K}(\alpha))\) . But the latter group operates simply-transitively on the set of framings of (G,T) whose first component is \(B_{0}\) .

PROPOSITION 19. The group \(\operatorname{Int}(\mathbf{G})\) operates simply-transitively on the set of framings of \(\mathbf{G}\) .

Let \(e = (\mathrm{T}, \mathrm{B}, (U_{\alpha}))\) and \(e' = (\mathrm{T}', \mathrm{B}', (U_{\alpha}')\) be two framings of \(G\) . There exist elements \(g\) in \(G\) such that \((\operatorname{Int} g)(\mathrm{T}) = \mathrm{T}'\) , and these elements form a single coset modulo \(N_G(T)\) . Thus, we can assume that \(T = T'\) , and we must prove that there exists a unique element of \(\operatorname{Int}_G(N_G(T))\) that transforms \(e\) to \(e'\) . By Chap. VI, §1, no. 5, Remark 4, there exists a unique element \(w\) of \(W(R)\) such that \(w(B) = B'\) . Since \(W(R)\) can be identified with \(N_G(T)/T\) , there exists \(n \in N_G(T)\) such that \(w = \operatorname{Int} n\) , and \(n\) is uniquely determined modulo \(T\) . Thus, we can assume that \(B = B'\) , and we must prove that there exists a unique element of \(\operatorname{Int}_{\mathrm{G}}(\mathrm{T})\) that transforms e to \(e'\) , which is simply Lemma 3.

COROLLARY. Let e be a framing of (G,T) and let E be the group of automorphisms of G that leave e stable. Then \(\operatorname{Aut}(G)\) is the semi-direct product of E by \(\operatorname{Int}(G)\) , and \(\operatorname{Aut}(G,T)\) is the semi-direct product of E by \(\operatorname{Int}(G)\cap\operatorname{Aut}(G,T)=\operatorname{Int}_{G}(\mathrm{N}_{\mathrm{G}}(\mathrm{T}))\) .

Indeed, every element of \(\operatorname{Aut}(G)\) transforms e into a framing of G. By Prop. 19, every coset of \(\operatorname{Aut}(G)\) modulo \(\operatorname{Int}(G)\) meets E in a single point, hence the first assertion. The second is proved in the same way.

Remark. Let G and \(G'\) be two connected compact Lie groups, and let \(e = (\mathrm{T}, \mathrm{B}, (U_{\alpha}))\) and \(e' = (\mathrm{T}', \mathrm{B}', (U_{\alpha}')\) be framings of G and \(G'\) , respectively. Let X be the set of isomorphisms from G to \(G'\) that take e to \(e'\) . The map \(f \mapsto \mathrm{D}^{*}(f)\) (resp. \(\mathrm{D}_{*}(f)\) ) is a bijection from X to the set of isomorphisms from \(\mathrm{D}^{*}(\mathrm{G}', \mathrm{T}')\) to \(\mathrm{D}^{*}(\mathrm{G}, \mathrm{T})\) (resp. \(\mathrm{D}_{*}(\mathrm{G}, \mathrm{T})\) to \(\mathrm{D}_{*}(\mathrm{G}', \mathrm{T}')\) ) that map \(B'\) to B (resp. B to \(B'\) ). Indeed, this follows immediately from Prop. 15 and Lemma 3.

\section{CONJUGACY CLASSES}

We retain the notations of §4.

\subsection{REGULAR ELEMENTS}

By Cor. 4 of Th. 2 of §2, no. 2, the regular elements \(g\) of \(G\) can be characterized by either of the following properties:

\begin{enumerate}
\def\labelenumi{\alph{enumi})}
\item
  The subalgebra of g fixed by Ad g is a Cartan subalgebra.
\item
  \(Z(g)_{0}\) is a maximal torus of G.
\end{enumerate}

The set of regular elements of \(\mathrm{G}\) is open and dense in \(\mathrm{G}\) .

In the remainder of this paragraph, we denote by \(G_{r}\) (resp. \(T_{r}\) ) the set of points of G (resp. T) that are regular in G. An element g of G belongs to \(T_{r}\) if and only if \(Z(g)_{0}\) is equal to T; every element of \(G_{r}\) is conjugate to an element of \(T_{r}\) (§2, no. 2, Th. 2).

An element t of T belongs to \(T_{r}\) if and only if \(t^{\alpha} \neq 1\) for every root \(\alpha \in \mathrm{R}(\mathrm{G}, \mathrm{T})\) ; consequently, \(T - T_{r}\) is the union of the subtori Ker \(\alpha\) for \(\alpha\) in \(\mathrm{R}(\mathrm{G}, \mathrm{T})\) .

PROPOSITION 1. Put \(n = \dim G\) . There exists a compact real-analytic manifold \(V\) of dimension \(n - 3\) and an analytic map \(\varphi : V \to G\) whose image is \(G - G_r\) .

Let \(\alpha\in\mathrm{R}(\mathrm{G},\mathrm{T})\) ; put \(\mathrm{V}_{\alpha}=(\mathrm{G}/\mathrm{Z}(\mathrm{Ker}\alpha))\times(\mathrm{Ker}\alpha)\) , and let \(\varphi_{\alpha}\) be the morphism from \(V_{\alpha}\) to G such that, for all \(g\in G\) and all \(t\in\operatorname{Ker}\alpha\) , we have \(\varphi_{\alpha}(\bar{g},t) = gtg^{-1}\) (we denote by \(\bar{g}\) the coset of \(g\) modulo \(\mathrm{Z}(\mathrm{Ker}\alpha)\) ). Then \(\mathrm{V}_{\alpha}\) is a compact real-analytic manifold of dimension

\[
\begin{array}{c} \dim V _ {\alpha} = \dim G - \dim Z (\operatorname{Ker} \alpha) + \dim \operatorname{Ker} \alpha \\ = n - (\dim T + 2) + (\dim T - 1) = n - 3 \end{array}
\]

(§4, no. 5, Th. 1); \(\varphi_{\alpha}\) is a morphism of real-analytic manifolds, and the image of \(\varphi_{\alpha}\) consists of the elements of G conjugate to an element of \(\mathrm{Ker} \alpha\) . It now suffices to take V to be the sum of the manifolds \(V_{\alpha}\) , and \(\varphi\) to be the morphism inducing \(\varphi_{\alpha}\) on each \(V_{\alpha}\) .

Remark. Call an element g of G very regular if \(Z(g)\) is a maximal torus of G. If \(g \in T\) , g is very regular if and only if \(w(g) \neq g\) for every non-identity element w of \(W_{G}(T)\) (§4, no. 7, Prop. 14). Thus, the set of very regular elements is a dense open subset of G (§2, no. 5, Cor. 2 of Prop. 5).

\subsection{CHAMBERS AND ALCOVES}

Denote by \(t_{r}\) the set of elements \(x \in t\) such that exp x is regular, in other words belongs to \(T_{r}\) . An element x of t belongs to \(t - t_{r}\) if and only if there exists a root \(\alpha \in \mathrm{R}(\mathrm{G}, \mathrm{T})\) such that \(\delta(\alpha)(x) \in 2\pi i\mathbf{Z}\) . For each root \(\alpha \in \mathrm{R}(\mathrm{G}, \mathrm{T})\) and each integer n, denote by \(H_{\alpha,n}\) the set of \(x \in t\) such that \(\delta(\alpha)(x) = 2\pi in\) . The \(H_{\alpha,n}\) are called the singular hyperplanes of t, and \(t - t_{r}\) is the union of the singular hyperplanes. The alcoves of t are the connected components of \(t_{r}\) , and the chambers are the connected components of the complement in t of the union of those singular hyperplanes that pass through the origin (that is, the \(H_{\alpha,0} = \operatorname{Ker}\delta(\alpha), \alpha \in \mathrm{R}(\mathrm{G}, \mathrm{T})\) ).

We have \(\varGamma(\mathrm{T})\subset\mathfrak{t}-\mathfrak{t}_{r}\) ; denote by N(G,T) the subgroup of \(\varGamma(\mathrm{T})\) generated by the nodal vectors (§4, no. 5); by Prop. 11 of §4, no. 6, the quotient \(\varGamma(\mathrm{T})/\mathrm{N}(\mathrm{G},\mathrm{T})\) can be identified with the fundamental group of G.

Finally, denote by W the Weyl group of G relative to T, considered as a group of automorphisms of T and of t, and denote by \(W_{a}\) (resp. \(W_{a}^{\prime}\) ) the group of automorphisms of the affine space t generated by W and the translations \(t_{\gamma}: x \mapsto x + \gamma\) for \(\gamma \in \mathrm{N}(\mathrm{G}, \mathrm{T})\) (resp. for \(\gamma \in \Gamma(\mathrm{T})\) ).

Let \(w \in W\) , \(\gamma \in \Gamma(T)\) , \(\alpha \in R(G, T)\) and \(n \in Z\) . We have:

\[
w (\mathrm{H} _ {\alpha , n}) = \mathrm{H} _ {w \alpha , n}, t _ {\gamma} (\mathrm{H} _ {\alpha , n}) = \mathrm{H} _ {\alpha , n + \langle \gamma , \alpha \rangle}.
\]

It follows that, for all chambers C and all \(w \in W\) , \(w(\mathrm{C})\) is a chamber and that for all alcoves A and all \(w \in W_{a}^{\prime}\) , \(w(\mathrm{A})\) is an alcove. Note that, when \(\mathrm{X}(\mathrm{T}) \otimes \mathbf{R}\) is identified with \(t^{*}\) via the isomorphism \((2\pi i)^{-1}\delta\) , the alcoves of t and the group \(W_{a}\) are the alcoves and the affine Weyl group associated to the root system \(\mathrm{R}(\mathrm{G}, \mathrm{T})\) (Chap. VI, §2, no. 1).

PROPOSITION 2. a) The group \(W_{a}\) (resp. \(W_{a}^{\prime}\) ) is the semi-direct product of W by N(G, T) (resp. \(\Gamma(T)\) ); the subgroup \(W_{a}\) of \(W_{a}^{\prime}\) is normal.

\begin{enumerate}
\def\labelenumi{\alph{enumi})}
\setcounter{enumi}{1}
\item
  The group W (resp. \(\mathrm{W}_a\) ) operates simply-transitively on the set of chambers (resp. alcoves).
\item
  Let C be a chamber and A an alcove. Then \(\overline{C}\) (resp. \(\overline{A}\) , resp. A) is a fundamental domain for the operation of W on t (resp. of \(W_{a}\) on t, resp. of \(W_{a}\) on \(t-t_{r}\) ). If \(x \in t_{r}\) and \(w \in W_{a}\) are such that \(w(x) = x\) , then w = Id.
\item
  For every chamber C, there exists a unique alcove A such that \(A \subset C\) and \(0 \in \overline{A}\) . For every alcove A, there exists a unique \(\gamma \in \mathrm{N}(\mathrm{G}, \mathrm{T})\) such that \(\gamma \in \overline{A}\) .
\end{enumerate}

If \(w \in W\) and \(\gamma \in \Gamma(T)\) , then \(wt_{\gamma}w^{-1} = t_{w(\gamma)}\) and \(wt_{\gamma}w^{-1}t_{\gamma}^{-1} = t_{w(\gamma)-\gamma}\) , with \(w(\gamma) - \gamma \in \mathrm{N}(\mathrm{G}, \mathrm{T})\) ; this immediately implies a). The rest of the proposition follows from Chap. VI, §1, no. 5 and §2, nos. 1 and 2.

COROLLARY 1. Let A be an alcove of t, \(\overline{A}\) its closure, and \(H_{A}\) the stabilizer of A in \(W_{a}^{\prime}\) .

\begin{enumerate}
\def\labelenumi{\alph{enumi})}
\item
  The group \(W_{a}^{\prime}\) is the semi-direct product of \(H_{A}\) by \(W_{a}\) .
\item
  The exponential map \(\overline{\mathrm{A}}\to \mathrm{T}\) and the canonical injection \(\mathrm{T}\rightarrow \mathrm{G}\) induce by passage to the quotients and to subsets homeomorphisms
\end{enumerate}

\[
\begin{array}{l} \overline {{\mathrm{A}}} / \mathrm {H_ {A}} \to \mathrm{T} / \mathrm{W} \to \mathrm{G} / \mathrm{Int} (\mathrm{G}) \\ \mathrm{A} / \mathrm {H_ {A}} \to \mathrm {T_ {r}} / \mathrm{W} \to \mathrm {G_ {r}} / \mathrm{Int} (\mathrm{G}). \end{array}
\]

Let \(w' \in W_a'\) ; then \(w'(A)\) is an alcove of \(t\) , and there exists (Prop. 2 b)) a unique element \(w\) of \(W_a\) such that \(w(A) = w'(A)\) , that is \(w^{-1}w' \in H_A\) . Since \(W_a\) is normal in \(W_a'\) , this proves \(a\) ).

The canonical injection of \(\overline{A}\) into t induces a continuous bijection \(\theta: \overline{A} \to t/W_{a}\) (Prop. 2 c)), and this is a homeomorphism since \(\overline{A}\) is compact. Since \(W_{a}\) is normal in \(W_{a}^{\prime}\) , the group \(H_{A}\) operates canonically on \(t/W_{a}\) (Algebra, Chap. I, §5, no. 4) and \(t/W_{a}^{\prime}\) can be identified with the quotient \((t/W_{a})/H_{A}\) ; the map \(\theta\) is compatible with the operations of \(H_{A}\) , hence induces by passage to the quotient a homeomorphism \(\overline{A}/H_{A} \to t/W_{a}^{\prime}\) . Further, \(\exp_{\Gamma}\) induces a homeomorphism from \(t/\Gamma(T)\) to T, hence also a homeomorphism from \(t/W_{a}^{\prime}\) to T/W. Assertion b) follows from that and Cor. 1 of Prop. 5 of §2, no. 4.

Remarks. 1) The group \(H_{A}\) can be identified naturally with \(\Gamma(\mathrm{T})/\mathrm{N}(\mathrm{G},\mathrm{T})\) , hence also with \(\pi_{1}(\mathrm{G})\) . Thus, it reduces to the identity element when G is simply-connected.

\begin{enumerate}
\def\labelenumi{\arabic{enumi})}
\setcounter{enumi}{1}
\item
  Let \(x \in \mathrm{A}\) ; then \(\exp x \in \mathrm{T}_r\) , so \(\mathrm{Z}(\exp x)_0 = \mathrm{T}\) . Further, \(\exp x\) is very regular (no. 1, Remark) if and only if \(w(x) \neq x\) for all \(w \in \mathrm{W}_a'\) distinct from the identity. By Cor. 1, this also means that \(h(x) \neq x\) for all \(h \in \mathrm{H}_{\mathrm{A}}\) distinct from the identity. In particular, if \(\mathrm{G}\) is simply-connected, then \(\mathrm{Z}_{\mathrm{G}}(t) = \mathrm{T}\) for all \(t \in \mathrm{T}_r\) , and every regular element of \(\mathrm{G}\) is very regular.
\item
  The special points of \(W_{a}\) (Chap. VI, §2, no. 2) are the elements x of t such that \(\delta(\alpha)(x) \in 2\pi i\mathbf{Z}\) for all \(\alpha \in \mathrm{R}(\mathrm{G}, \mathrm{T})\) (loc. cit., Prop. 3), that is such that \(\exp x \in \mathrm{C}(\mathrm{G})\) (§4, no. 4, Prop. 8). For such an element x we have \(wx - x \in \mathrm{N}(\mathrm{G}, \mathrm{T})\) for all \(w \in W\) (Chap. VI, §1, no. 9, Prop. 27), so the stabilizers of x in \(W_{a}\) and \(W_{a}^{\prime}\) coincide. Let S be the set of special points of \(\overline{A}\) ; it follows from the preceding and Cor. 1 that the group \(H_{A}\) operates freely on S, and that the exponential map induces a bijection from \(S/H_{A}\) to \(C(G)\) .
\end{enumerate}

COROLLARY 2. Let C be a chamber of t and \(\overline{C}\) its closure. The canonical injections \(\overline{C} \rightarrow t \rightarrow g\) induce by passage to the quotient homeomorphisms

\[
\overline {{\mathrm{C}}} \rightarrow \mathfrak {t} / \mathrm{W} \rightarrow \mathfrak {g} / \operatorname{Ad} (\mathrm{G}).
\]

The canonical maps \(\overline{C} \rightarrow t\) and \(t \rightarrow t/W\) are proper (General Topology, Chap. III, §4, no. 1, Prop. 2 c)). The map \(\overline{C} \rightarrow t/W\) is continuous, proper and bijective (Prop. 2 c)); thus, it is a homeomorphism, hence the corollary in view of the Cor. of Prop. 6 of §2, no. 5.

Remark 4. Denote by \(\mathfrak{g}_{\mathrm{reg}}\) the set of regular elements of \(\mathfrak{g}\) (Chap. VII, §2, no. 2, Def. 2) and put \(\mathfrak{t}_{\mathrm{reg}} = \mathfrak{t} \cap \mathfrak{g}_{\mathrm{reg}}\) . For \(x \in \mathfrak{t}\) , we have

\[
\det (\mathrm{X} - \operatorname{ad} _ {\mathfrak {g}} x) = \mathrm{X} ^ {\dim t} \prod_ {\alpha \in \mathrm{R(G,T)}} (\mathrm{X} - \delta (\alpha) (x)),
\]

and hence \(t_{reg}\) is the set of elements x of t such that \(\delta(\alpha)(x) \neq 0\) for all \(\alpha \in \mathrm{R}(\mathrm{G}, \mathrm{T})\) , that is the union of the chambers of t (so \(t_r \subset t_{reg}\) ). Consequently \(\overline{C} \cap t_{reg} = C\) , so we have homeomorphisms

\[
\mathrm{C} \rightarrow \mathfrak {t} _ {\text { reg }} / \mathrm{W} \rightarrow \mathfrak {g} _ {\text { reg }} / \mathrm{Ad} (\mathrm{G}).
\]

COROLLARY 3. Assume that G is simply-connected; let g be a regular element of G. There exist a maximal torus S of G, and an alcove A of L(S), both uniquely determined, such that \(g \in \exp(A)\) and \(0 \in \overline{A}\) .

We can assume that \(g\) belongs to \(\mathrm{T}_r\) (§2, no. 2, Th. 2). Let \(x\) be an element of \(\mathfrak{t}_r\) such that \(\exp x = g\) , and let \(\mathrm{A}'\) be the alcove of \(\mathfrak{t}\) containing \(x\) . The alcoves \(\mathrm{A}\) of \(\mathfrak{t}\) such that \(g \in \exp(\mathrm{A})\) are the alcoves \(\mathrm{A}' - \gamma\) for \(\gamma \in \Gamma(\mathrm{T})\) ; thus, the assertion follows from Prop. 2 d).

\subsection{AUTOMORPHISMS AND REGULAR ELEMENTS}

Lemma 1. Let \(u\) be an automorphism of \(G\) , and \(H\) the set of its fixed points. a) \(H\) is a closed subgroup of \(G\) .

\begin{enumerate}
\def\labelenumi{\alph{enumi})}
\setcounter{enumi}{1}
\tightlist
\item
  If \(\mathrm{H}_0\) is central in G, then G is commutative (so \(\mathrm{G} = \mathrm{T}\) ).
\end{enumerate}

Assertion \(a\) ) is clear. To prove \(b\) ), we can replace \(\mathrm{G}\) by \(\mathrm{D}(\mathrm{G})\) ( \(\S 1\) , Cor. 1 of Prop. 4), and hence can assume that \(\mathrm{G}\) is semi-simple. Then, if \(\mathrm{H}_0\) is central in \(\mathrm{G}\) , we have \(\mathrm{L}(\mathrm{H}) = \{0\}\) , so the endomorphism \(\mathrm{L}(u) - \mathrm{Id}\) of \(\mathfrak{g}\) is bijective. Let \(f\) be the endomorphism of the manifold \(\mathrm{G}\) defined by \(f(g) = u(g)^{-1}g\) for \(g \in \mathrm{G}\) ; it is étale, for if \(g \in \mathrm{G}\) and \(x \in \mathfrak{g}\) , we have \(\mathrm{T}(f)(xg) = u(g)^{-1}(x - \mathrm{L}(u)(x))g\) , so the tangent map of \(f\) at \(g\) is bijective. It follows that the image of \(f\) is open and compact, hence coincides with \(\mathrm{G}\) since \(\mathrm{G}\) is connected. Now let \(\mathrm{E}\) be a framing of G (§4, no. 10, Def. 3) and \(u(\mathrm{E})\) its image under \(u\) . By Prop. 19 of loc. cit., there exists an element \(h\) of G such that \((\operatorname{Int} h)(\mathrm{E}) = u(\mathrm{E})\) . Let \(g \in \mathrm{G}\) be such that \(h = f(g) = u(g)^{-1}g\) ; then

\[
u \circ \operatorname{Int} g = (\operatorname{Int} u (g)) \circ u = \operatorname{Int} g \circ (\operatorname{Int} h) ^ {- 1} \circ u,
\]

so the framing \((\operatorname{Int} g)(\operatorname{E})\) is stable under u. If \((\operatorname{Int} g)(\operatorname{E}) = (\operatorname{T}_{1}, \operatorname{B}, (U_{\alpha})_{\alpha \in \operatorname{B}})\) , then \(\sum U_{\alpha} \in \operatorname{L}(\operatorname{H})\) ; since \(\operatorname{L}(\operatorname{H}) = \{0\}\) , this implies that \(B = \varnothing\) , so \(G = T_{1}\) , and G is commutative.

Lemma 2. Let x be an element of T and S a subtorus of T. If the identity component of \(Z(x) \cap Z(S)\) reduces to T, there exists an element s of S such that xs is regular.

For all \(\alpha\in\mathrm{R}(\mathrm{G},\mathrm{T})\) , let \(S_{\alpha}\) be the submanifold of S consisting of the elements s of S such that \((xs)^{\alpha}=1\) . If there is no element s of S such that xs is regular, S is the union of the submanifolds \(S_{\alpha}\) , hence is equal to one of them. Hence there exists \(\alpha\) in \(\mathrm{R}(\mathrm{G},\mathrm{T})\) such that \((xs)^{\alpha}=1\) for all \(s\in S\) ; but this implies that \(x^{\alpha}=1\) and \(\alpha|S=1\) , so \(\mathrm{Z}(x)\cap\mathrm{Z}(\mathrm{S})\supset\mathrm{Z}(\mathrm{Ker}\alpha)\) , hence the lemma.

Lemma 3. Assume that G is simply-connected. Let C be a chamber of t, and u an automorphism of G such that T and C are stable under u. Then the set of points of T fixed by u is connected.

Since G is simply-connected, \(\varGamma(\mathrm{T})\) is generated by the nodal vectors \(K_{\alpha}\) ( \(\alpha \in \mathrm{R}(\mathrm{G}, \mathrm{T})\) ), and hence has a basis consisting of the family of the \(K_{\alpha}\) for which \(\alpha\) belongs to the basis \(\mathrm{B}(\mathrm{C})\) defined by C (Chap. VI, §1, no. 10). Thus, it suffices to prove that, if \(\varphi\) is an automorphism of the torus T leaving stable a basis of \(\varGamma(\mathrm{T})\) , the set of fixed points of \(\varphi\) is connected. Decomposing this basis into the disjoint union of the orbits of the group generated by \(\varphi\) , we are reduced to the case in which \(T = U^{n}\) and \(\varphi\) is the automorphism \((z_{1}, \ldots, z_{n}) \mapsto (z_{2}, \ldots, z_{n}, z_{1})\) ; in this case the fixed points of \(\varphi\) are the points \((z, z, \ldots, z)\) for \(z \in U\) , which form a connected subgroup of T.

PROPOSITION 3. Let \(u\) be an automorphism of \(G\) , and let \(x\) be a point of \(G\) fixed by \(u\) .

\begin{enumerate}
\def\labelenumi{\alph{enumi})}
\item
  There exists an element \(a\) of \(\mathfrak{g}\) , fixed by \(\mathrm{L}(u)\) and by \(\operatorname{Ad} x\) , such that \(x \exp a\) is regular.
\item
  There exists a regular element \(g\) of \(G\) fixed by \(u\) and commuting with \(x\) .
\end{enumerate}

Let H be the group of fixed points of u, S a maximal torus of \(\mathrm{Z}(x)\cap\mathrm{H}\) , and K the identity component of \(\mathrm{Z}(\mathrm{S})\cap\mathrm{Z}(x)\) . This is a connected closed subgroup of G; further, by Cor. 5 of Th. 2 of §2, no. 2, there exist maximal tori of G containing S and x, so K is of maximal rank and contains S and x. On the other hand, K is stable under u since S and x are; denote by V the set of fixed points of u in K. Then

\[
\mathrm{S} \subset \mathrm{V} _ {0} \subset \mathrm{K} \cap \mathrm{H} \subset \mathrm{Z} (\mathrm{S}) \cap \mathrm{Z} (x) \cap \mathrm{H},
\]

so \(V_{0}\) is contained in the centralizer of S in \((\mathrm{Z}(x)\cap\mathrm{H})_{0}\) ; but the latter reduces to S (loc. cit., Cor. 6), hence finally \(V_{0}=S\) . Lemma 1 now implies that K is commutative, hence is a maximal torus of G (since it is connected and of maximal rank). It contains S and x, and is equal to the identity component of \(\mathrm{Z}(\mathrm{S})\cap\mathrm{Z}(x)\) ; assertion a) now follows from Lemma 2. We deduce b) by taking \(g=x\exp a\) .

COROLLARY. Let \(\mathfrak{s}\) be a compact Lie algebra, and let \(\varphi\) be an automorphism of \(\mathfrak{s}\) . There exists a regular element of \(\mathfrak{s}\) fixed by \(\varphi\) .

Replacing \(\mathfrak{s}\) by \(\mathcal{D}\mathfrak{s}\) , we can assume that \(\mathfrak{s}\) is semi-simple. Let \(S\) be a compact simply-connected Lie group with Lie algebra \(\mathfrak{s}\) , and let \(u\) be the automorphism of \(S\) such that \(L(u) = \varphi\) . Prop. 3 implies the existence of an element \(a\) of \(\mathfrak{s}\) , fixed by \(\varphi\) , such that \(\exp a\) is regular in \(S\) ; in particular, \(a\) is regular in \(\mathfrak{s}\) (no. 2, remark 4).

THEOREM 1. Let \(u\) be an automorphism of a connected compact Lie group \(G\) .

\begin{enumerate}
\def\labelenumi{\alph{enumi})}
\item
  The identity component of the group of fixed points of \(u\) contains a regular element of \(G\) .
\item
  There exists a maximal torus K of G and a chamber of L(K) that are stable under u.
\item
  If G is simply-connected, the set of fixed points of u is connected.
\end{enumerate}

Assertion a) is the particular case x = e of Prop. 3. We assume now that G is simply-connected and prove b) and c). Let x be an element of G fixed by u, and let g be a regular element of G, fixed by u and commuting with x (Prop. 3). The centralizer K of g is a maximal torus of G (no. 2, Remark 2), stable under u, and containing x and g. By Cor. 3 of Prop. 2 of no. 2, there exists a unique alcove A of L(K) such that \(g \in \exp A\) and \(0 \in \overline{A}\) ; since g is fixed by u, L(u) leaves A, and hence also the chamber of L(K) containing A, stable. This proves b); further, the set of points of K fixed by u is connected (Lemma 3) and contains x and e, hence c) (General Topology, Chap. I, §11, no. 1, Prop. 2).

It remains to prove \(b)\) in the general case. Now, if \(\tilde{\mathrm{D}}(\mathrm{G})\) is the universal covering of \(\mathrm{D}(\mathrm{G})\) , and if \(f: \tilde{\mathrm{D}}(\mathrm{G}) \to \mathrm{G}\) is the canonical morphism, there exists an automorphism \(\tilde{u}\) of \(\tilde{\mathrm{D}}(\mathrm{G})\) such that \(f \circ \tilde{u} = u \circ f\) . If \(\tilde{\mathrm{K}}\) is a maximal torus of \(\tilde{\mathrm{D}}(\mathrm{G})\) and \(\tilde{\mathrm{C}}\) a chamber of \(\mathrm{L}(\tilde{\mathrm{K}})\) , stable under \(\tilde{u}\) (this exists by what has already been proved), there exists (§2, no. 3, Remark 2) a unique maximal torus \(\mathrm{K}\) of \(\mathrm{G}\) and a unique chamber \(\mathrm{C}\) of \(\mathrm{L}(\mathrm{K})\) such that \(\tilde{\mathrm{K}} = f^{-1}(\mathrm{K})\) and \(\tilde{\mathrm{C}} = \mathrm{L}(f)^{-1}(\mathrm{C})\) , and we see immediately that \(\mathrm{K}\) and \(\mathrm{C}\) are stable under \(u\) , hence assertion \(b)\) in the general case.

COROLLARY 1. Assume that the Z-module \(\pi_{1}(G)\) is torsion-free.

\begin{enumerate}
\def\labelenumi{\alph{enumi})}
\item
  The centralizer of every element of G is connected.
\item
  Any two commuting elements of G belong to the same maximal torus.
\end{enumerate}

By Cor. 3 of Prop. 11 of §4, no. 6, D(G) is simply-connected. We have G = C(G) \(_0\) .D(G); let \(x \in\) G; write \(x = uv\) , with \(u \in\) C(G) \(_0\) and \(v \in\) D(G). Then Z(x) = C(G) \(_0\) .Z \(_{\text{D(G)}}\) (v). By Th. 1 c), Z \(_{\text{D(G)}}\) (v) is connected, so Z(x) is connected, hence a). Thus, by Cor. 3 of Th. 2 of §2, no. 2, Z(x) is the union of the maximal tori of G containing \(x\) , hence b).

COROLLARY 2. Let \(\Gamma\) be a compact subgroup of \(\operatorname{Aut}(\mathbf{G})\) with the following property:

(*) There exist elements \(u_{1}, \ldots, u_{n}\) of \(\Gamma\) such that, for all \(i\) , the closure \(\Gamma_{i}\) of the subgroup of \(\Gamma\) generated by \(u_{1}, \ldots, u_{i}\) is a normal subgroup of \(\Gamma\) , and \(\Gamma_{n} = \Gamma\) .

Then, there exists a maximal torus of G stable under the operation of \(\Gamma\) .

We argue by induction on the dimension of G. Clearly, we can assume that \(u_{1} \neq Id\) ; then the subgroup H of fixed points of \(u_{1}\) is distinct from G, and is stable under the operation of \(\varGamma\) . Moreover, since \(\varGamma\) is compact, the image of \(\varGamma\) in \(\mathrm{Aut}(\mathrm{H}_{0})\) is a quotient of \(\varGamma\) , hence it also satisfies condition (*). By the induction hypothesis, there exists a maximal torus S of H stable under \(\varGamma\) . The centralizer K of S in G is connected ( \(\S2\) , no. 2, Cor. 5) and stable under \(\varGamma\) ; this is a maximal torus of G, since \(H_{0}\) contains a regular element of G (Th. 1 a)) which is conjugate to an element of S (loc. cit., Cor. 4).

COROLLARY 3. Let H be a Lie group and \(\Gamma\) a compact subgroup of H. Assume that \(\mathrm{H}_0\) is compact and that \(\Gamma\) satisfies condition (*) of Cor. 2. Then there exists a maximal torus T of \(\mathrm{H}_0\) such that \(\Gamma \subset \mathrm{N}_{\mathrm{H}}(\mathrm{T})\) .

COROLLARY 4. Every nilpotent subgroup of a compact Lie group is contained in the normalizer of a maximal torus.

Let H be a compact Lie group, N a nilpotent subgroup of H. Then the closure \(\varGamma\) of N is also a nilpotent group (Chap. III, §9, no. 1, Cor. 2 of Prop. 1), and it suffices, in view of Cor. 3, to prove that \(\varGamma\) satisfies condition (*). Now \(\varGamma_{0}\) is a connected compact nilpotent Lie group, hence is a torus (§1, no. 4, Cor. 1 of Prop. 4), and there exists an element \(u_{1}\) of \(\varGamma\) generating a dense subgroup of \(\varGamma_{0}\) (General Topology, Chap. VII, §1, no. 3, text preceding Prop. 8). The finite group \(\varGamma/\varGamma_{0}\) is nilpotent and there exist \(\tilde{u}_{2},\ldots,\tilde{u}_{n}\in\varGamma/\varGamma_{0}\) generating \(\varGamma/\varGamma_{0}\) and such that the subgroup of \(\varGamma/\varGamma_{0}\) generated by \((\tilde{u}_{2},\ldots,\tilde{u}_{r})\) is normal for \(r=2,\ldots,n\) (Algebra, Chap. I, §6, no. 5, Th. 1 and no. 7, Th. 4). Then, if \(u_{2},\ldots,u_{n}\) are representatives of \(\tilde{u}_{2},\ldots,\tilde{u}_{n}\) in \(\varGamma\) , the sequence \((u_{1},\ldots,u_{n})\) has the required properties.

Example. Take \(\mathbf{G} = \mathbf{U}(n,\mathbf{C})\) . We shall see later that the subgroup of diagonal matrices in \(\mathbf{G}\) is a maximal torus of \(\mathbf{G}\) and that its normalizer is the set of monomial matrices (Algebra, Chap. II, §10, no. 7, Example II) in \(\mathbf{G}\) .

We conclude that, if \(\Phi\) is a separating positive hermitian form on a finite dimensional complex vector space V and \(\Gamma\) is a nilpotent subgroup of \(\mathbf{U}(\Phi)\) , there exists a basis of V for which the matrices of the elements of \(\Gamma\) are monomial (``Blichtfeldt's theorem'').

\subsection{\texorpdfstring{THE MAPS \((\mathbf{G} / \mathbf{T})\times \mathbf{T}\to \mathbf{G}\) AND \((\mathbf{G} / \mathbf{T})\times \mathbf{A}\rightarrow \mathbf{G}_r\)}{THE MAPS (\textbackslash mathbf\{G\} / \textbackslash mathbf\{T\})\textbackslash times \textbackslash mathbf\{T\}\textbackslash to \textbackslash mathbf\{G\} AND (\textbackslash mathbf\{G\} / \textbackslash mathbf\{T\})\textbackslash times \textbackslash mathbf\{A\}\textbackslash rightarrow \textbackslash mathbf\{G\}\_r}}

The map \((g,t)\mapsto gtg^{-1}\) from G × T to G induces by passage to the quotient a morphism of analytic manifolds

\[
f: (\mathrm{G} / \mathrm{T}) \times \mathrm{T} \to \mathrm{G},
\]

which is surjective (§2, no. 2, Th. 2). By restriction, \(f\) induces a surjective morphism

\[
f _ {r}: (\mathrm{G} / \mathrm{T}) \times \mathrm{T} _ {r} \rightarrow \mathrm{G} _ {r}.
\]

By composition with \(Id_{G/T} \times exp_{T}\) , we obtain morphisms, also surjective,

\[
\begin{array}{c} \varphi : (\mathrm{G/T}) \times \mathfrak {t} \to \mathrm{G}, \\ \varphi_ {r}: (\mathrm{G/T}) \times \mathfrak {t} _ {r} \to \mathrm{G} _ {r}; \end{array}
\]

finally, if A is an alcove of t, \(\varphi_{r}\) induces a surjective morphism

\[
\varphi_ {\mathrm{A}}: (\mathrm{G/T}) \times \mathrm{A} \to \mathrm{G} _ {r}.
\]

We define a right operation of W on G/T as follows: let \(w \in W\) and \(u \in G/T\) ; lift w to an element n of \(\mathrm{N}_{\mathrm{G}}(\mathrm{T})\) and u to an element g of G. Then the image of gn in G/T does not depend on the choice of n and g; we denote it by u.w.

For this operation, W operates freely on G/T: indeed, with the preceding notations, assume that u.w = u; then \(gn \in gT\) , so \(n \in T\) and w = 1.

We define a right operation of W on \((\mathrm{G}/\mathrm{T}) \times \mathrm{T}\) by

\[
(u, t). w = (u. w, w ^ {- 1} (t)), \quad u \in \mathrm{G} / \mathrm{T}, t \in \mathrm{T}, w \in \mathrm{W}
\]

and a right operation of \(W_{a}^{\prime}\) on (G/T) × t by

\[
(u, x). \omega = (u. \overline {{\omega}}, \omega^ {- 1} (x)), \quad u \in \mathrm{G/T}, x \in \mathfrak {t}, \omega \in \mathrm{W} _ {a} ^ {\prime},
\]

where \(\overline{\omega}\) is the image of \(\omega\) in the quotient \(W_{a}^{\prime}/\Gamma(T)=W\) .

If A is an alcove of t, and if \(H_{A}\) is the subgroup of \(W_{a}^{\prime}\) that stabilizes A, we obtain by restriction an operation of \(H_{A}\) on \((\mathrm{G}/\mathrm{T}) \times \mathrm{A}\) .

These different operations are compatible with the morphisms \(f, \varphi\) and \(\varphi_{\mathrm{A}}\) : for \(u \in \mathrm{G} / \mathrm{T}, t \in \mathrm{T}, x \in \mathfrak{t}, y \in \mathrm{A}, w \in \mathrm{W}, \omega \in \mathrm{W}_a', h \in \mathrm{H}_{\mathrm{A}}\) , we have

\[
f ((u, t). w) = f (u, t), \varphi ((u, x). \omega) = \varphi (u, x), \varphi_ {\mathrm{A}} ((u, y). h) = \varphi_ {\mathrm{A}} (u, y).
\]

Lemma 4. Let \(g \in \mathrm{G}\) , \(t \in \mathrm{T}\) , and let \(\bar{g}\) be the image of \(g\) in \(\mathrm{G} / \mathrm{T}\) . Identify the tangent space of \(\mathrm{G} / \mathrm{T}\) (resp. T, resp. G) at \(\bar{g}\) (resp. t, resp. \(gtg^{-1}\) ) with \(\mathfrak{g} / \mathfrak{t}\) (resp. \(\mathfrak{t}\) , resp. \(\mathfrak{g}\) ) by means of the left translation \(\gamma(g)\) by \(g\) (resp. \(t\) , resp. \(gtg^{-1}\) ). The tangent linear map of \(f\) at \((\bar{g}, t)\) is then identified with the linear map \(f': (\mathfrak{g} / \mathfrak{t}) \times \mathfrak{t} \to \mathfrak{g}\) defined as follows: if \(z \in \mathfrak{g}, x \in \mathfrak{t}\) , and if \(\bar{z}\) denotes the image of \(z\) in \(\mathfrak{g} / \mathfrak{t}\) , then

\[
f ^ {\prime} (\bar {z}, x) = (\mathrm{Ad} g t ^ {- 1}) (z - (\mathrm{Ad} t) z + x).
\]

Let \(\mathrm{F}\) be the map from \(\mathrm{G} \times \mathrm{T}\) to \(\mathrm{T}\) such that \(\mathrm{F}(g,t) = gtg^{-1}\) . Since \(\mathrm{F} \circ (\gamma(g), \mathrm{Id}_{\mathrm{T}}) = \operatorname{Int} g \circ \mathrm{F}\) , we have \(\mathrm{T}_{(g,t)}(\mathrm{F})(gz, tx) = \mathrm{T}_t(\operatorname{Int} g) \circ \mathrm{T}_{(e,t)}(\mathrm{F})(z, tx)\) ; by Chap. III, §3, no. 12, Prop. 46,

\[
\mathrm{T} _ {(e, t)} (\mathrm{F}) (z, t x) = t ((\mathrm{Ad} t ^ {- 1}) z - z) + t x = t ((\mathrm{Ad} t ^ {- 1}) (z - (\mathrm{Ad} t) z + x))
\]

and consequently

\[
\mathrm{T} _ {(g, t)} (\mathrm{F}) (g z, t x) = g t g ^ {- 1} ((\mathrm{Ad}   g t ^ {- 1}) (z - (\mathrm{Ad}   t) z + x)).
\]

The lemma follows immediately from this formula by passage to the quotient.

PROPOSITION 4. a) Let \(g \in \mathbf{G}, t \in \mathrm{T}, x \in \mathfrak{t}\) , and let \(\bar{g}\) be the image of \(g\) in \(\mathrm{G} / \mathrm{T}\) . The following conditions are equivalent:

\begin{enumerate}
\def\labelenumi{(\roman{enumi})}
\tightlist
\item
  \(t \in \mathrm{T}_r\) (resp. \(x \in \mathfrak{t}_r\) ).
\end{enumerate}

(i bis) The element \(f(\bar{g}, t)\) (resp. \(\varphi(\bar{g}, x)\) ) is regular in G.

\begin{enumerate}
\def\labelenumi{(\roman{enumi})}
\setcounter{enumi}{1}
\tightlist
\item
  The map \(f\) (resp. \(\varphi\) ) is a submersion at the point \((\bar{g}, t)\) (resp. \((\bar{g}, x)\) ).
\end{enumerate}

(ii bis) The map \(f\) (resp. \(\varphi\) ) is étale at the point \((\bar{g}, t)\) (resp. \((\bar{g}, x)\) ).

\begin{enumerate}
\def\labelenumi{\alph{enumi})}
\setcounter{enumi}{1}
\item
  The map \(f_{r}\) (resp. \(\varphi_{r}\) , resp. \(\varphi_{A}\) ) makes \((\mathrm{G}/\mathrm{T}) \times \mathrm{T}_{r}\) (resp. \((\mathrm{G}/\mathrm{T}) \times \mathfrak{t}_{r}\) , resp. \((\mathrm{G}/\mathrm{T}) \times \mathrm{A}\) ) a principal covering of \(G_{r}\) with group W (resp. \(W_{a}^{\prime}\) , resp. \(H_{A}\) ).
\item
  The equivalence of (i) and (i bis) is clear; that of (ii) and (ii bis) follows from the relations \(\dim((\mathrm{G}/\mathrm{T})\times\mathrm{T})=\dim((\mathrm{G}/\mathrm{T})\times\mathfrak{t})=\dim(\mathrm{G})\) . By Lemma 4, f is a submersion at the point \((\bar{g},t)\) if and only if \(g=t+Im(Ad t-Id)\) , which means that t is regular. Finally, since \(\varphi=f\circ(Id_{G/T}\times\exp_{T})\) , \(\varphi\) is étale at the point \((\bar{g},x)\) if and only if f is étale at the point \((\bar{g},\exp x)\) , which by the preceding means that x belongs to \(t_{r}\) .
\item
  The morphisms \(f_{r}, \varphi_{r}, \varphi_{A}\) are thus étale. On the other hand, W operates freely on G/T, and a fortiori on \((\mathrm{G}/\mathrm{T}) \times \mathrm{T}\) . Let \(g, g'\) in G and \(t, t'\) in \(T_{r}\) be such that \(f(\bar{g}, t) = f(\bar{g}', t')\) ; then \(\operatorname{Int} g^{-1} g'\) maps \(t'\) to t, and hence normalizes T, since \(\mathrm{T} = \mathrm{Z}(t)_{0} = \mathrm{Z}(t')_{0}\) , and the class w of \(g^{-1} g'\) in W maps \((\bar{g}, t)\) to \((\bar{g}', t')\) . It follows that \(f_{r}\) is a principal covering with group W; this immediately implies that \(\varphi_{r}\) is a principal covering with group \(W_{a}'\) , and hence by restriction to the connected component \((\mathrm{G}/\mathrm{T}) \times \mathrm{A}\) of \((\mathrm{G}/\mathrm{T}) \times \mathbf{t}_{r}\) , that \(\varphi_{A}\) is a principal covering with group \(H_{A}\) .
\end{enumerate}

Remarks. 1) By Prop. 3 of §2, no. 4, the manifold \((\mathrm{G}/\mathrm{T}) \times \mathrm{A}\) is simply-connected. It follows that \(\varphi_{A}\) is a universal covering of \(G_{r}\) ; since \(\pi_{1}(G_{r})\) is canonically isomorphic to \(\pi_{1}(G)\) (no. 1, Prop. 1 and General Topology,

Chap. XI, in preparation), we recover the fact that \(\pi_{1}(\mathrm{G})\) can be identified with \(H_{A}\) (that is with \(\Gamma(\mathrm{T})/\mathrm{N}(\mathrm{G},\mathrm{T})\) ).

\begin{enumerate}
\def\labelenumi{\arabic{enumi})}
\setcounter{enumi}{1}
\item
  The restriction of \(\varphi_{A}\) to \(W \times A \subset (G/T) \times A\) makes \(W \times A\) a principal covering of \(T_{r}\) with group \(H_{A}\) . We thus recover Cor. 1 of Prop. 2 of no. 2.
\item
  Denote by \(\mathfrak{g}_r\) the inverse image of \(\mathrm{G}_r\) under the exponential map and by \(\varepsilon : \mathfrak{g}_r \to \mathrm{G}_r\) the map induced by \(\exp_{\mathrm{G}}\) . The map \((g, x) \mapsto (\operatorname{Ad} g)(x)\) from \(\mathrm{G} \times \mathfrak{t}_r\) to \(\mathfrak{g}_r\) defines by passage to the quotient a map \(\psi_r : (\mathrm{G} / \mathrm{T}) \times \mathfrak{t}_r \to \mathfrak{g}_r\) . We have \(\varepsilon \circ \psi_r = \varphi_r\) . Let \(w \in \mathrm{W}, \gamma \in \Gamma(\mathrm{T})\) and \(\omega \in \mathrm{W}_a'\) be such that \(\omega(z) = w(z) + \gamma\) for all \(z \in \mathfrak{t}\) ; then \(\psi_r((\bar{g}, x)\omega) = \psi_r(\bar{g}, x) - (\operatorname{Ad} g)(\gamma)\) for \(g \in \mathrm{G}, x \in \mathfrak{t}_r\) , so \(\psi_r((\bar{g}, x)\omega) = \psi_r(\bar{g}, x)\) if and only if \(\gamma = 0\) . It follows (cf.~General Topology, Chap. XI, in preparation) that \(\psi_r\) is a principal covering of \(\mathfrak{g}_r\) with group W, and that \(\varepsilon : \mathfrak{g}_r \to \mathrm{G}_r\) is a covering associated to the principal covering \(\varphi_r\) , with fibre isomorphic to the \(\mathrm{W}_a'\) -set \(\mathrm{W}_a'/\mathrm{W}\) .
\end{enumerate}

\section{INTEGRATION ON COMPACT LIE GROUPS}

We retain the notations of §4; put \(w(\mathrm{G}) = \mathrm{Card}(\mathrm{W}_{\mathrm{G}}(\mathrm{T}))\) . Denote by dg (resp. dt) the Haar measure on G (resp. T) with total mass 1, and by n (resp. r) the dimension of G (resp. T).

\subsection{PRODUCT OF ALTERNATING MULTILINEAR FORMS}

Let A be a commutative ring and M an A-module. For each integer \(r \geq 0\) , denote by \(\operatorname{Alt}^{r}(\mathbf{M})\) the A-module of alternating r-linear forms on M; it can be identified with the dual of the A-module \(\Lambda^{r}(\mathbf{M})\) (Algebra, Chap. III, §7, no. 4, Prop. 7). Let \(u \in \operatorname{Alt}^{s}(\mathbf{M})\) and \(v \in \operatorname{Alt}^{r}(\mathbf{M})\) ; recall (Algebra, Chap. III, §11, no. 2, Example 3) that the alternating product of u and v is the element \(u \wedge v \in \operatorname{Alt}^{s+r}(\mathbf{M})\) defined by

\[
(u \wedge v) (x _ {1}, \dots , x _ {s + r}) = \sum_ {\sigma \in \mathfrak {S} _ {s, r}} \varepsilon_ {\sigma} u (x _ {\sigma (1)}, \dots , x _ {\sigma (s)}) v (x _ {\sigma (s + 1)}, \dots , x _ {\sigma (s + r)}),
\]

where \(S_{s,r}\) is the subset of the symmetric group \(S_{s+r}\) consisting of the permutations whose restrictions to \((1,s)\) and \((s+1,s+r)\) are increasing.

Now let

\[
0 \longrightarrow \mathrm{M} ^ {\prime} \stackrel {i} {\longrightarrow} \mathrm{M} \stackrel {p} {\longrightarrow} \mathrm{M} ^ {\prime \prime} \longrightarrow 0
\]

be an exact sequence of free A-modules, of ranks \(r, r + s\) and \(s\) , respectively.

Lemma 1. There exists an A-bilinear map from \(\mathrm{Alt}^s (\mathbf{M}'')\times \mathrm{Alt}^r (\mathbf{M}')\) to \(\mathrm{Alt}^{s + r}(\mathbf{M})\) , denoted by \((u,v)\mapsto u\cap v\) , and characterized by either of the following two properties:

\begin{enumerate}
\def\labelenumi{\alph{enumi})}
\item
  Denote by \(u_{1} \in \operatorname{Alt}^{s}(\operatorname{M})\) the form \((x_{1}, \ldots, x_{s}) \mapsto u(p(x_{1}), \ldots, p(x_{s}))\) , and let \(v_{1} \in \operatorname{Alt}^{r}(\operatorname{M})\) be a form such that \(v_{1}(i(x_{1}^{\prime}), \ldots, i(x_{r}^{\prime})) = v(x_{1}^{\prime}, \ldots, x_{r}^{\prime})\) for \(x_{1}^{\prime}, \ldots, x_{r}^{\prime}\) in \(M^{\prime}\) ; then \(u \cap v = u_{1} \wedge v_{1}\) .
\item
  For all \(x_{1},\ldots,x_{s}\) in M and \(x_{1}^{\prime},\ldots,x_{r}^{\prime}\) in \(M^{\prime}\) ,
\end{enumerate}

\[
(u \cap v) (x _ {1}, \dots , x _ {s}, i (x _ {1} ^ {\prime}), \dots , i (x _ {r} ^ {\prime})) = u (p (x _ {1}), \dots , p (x _ {s})) v (x _ {1} ^ {\prime}, \dots , x _ {r} ^ {\prime}).\tag{1}
\]

The map \(\varphi : \mathrm{Alt}^s (\mathrm{M}'')\otimes_{\mathrm{A}}\mathrm{Alt}^r (\mathrm{M}')\to \mathrm{Alt}^{s + r}(\mathrm{M})\) such that \(\varphi (u\otimes v) = u\cap v\) is an isomorphism of free A-modules of rank one.

The existence of a form \(v_{1}\) satisfying condition \(a\) ) follows from the fact that \(\Lambda^r (i)\) induces an isomorphism from \(\Lambda^r (\mathrm{M}')\) to a direct factor submodule of \(\Lambda^r (\mathrm{M})\) (Algebra, Chap. III, §7, no. 2). Let \(v_{1}\) be such a form; put \(u\cap v = u_{1}\wedge v_{1}\) . Formula (1) is then satisfied, since if we put \(i(x_k') = x_{s + k}\) for \(1\leq k\leq r\) , the only element \(\sigma\) of \(\mathfrak{S}_{s,r}\) such that \(p(x_{\sigma (i)})\neq 0\) for \(1\leq i\leq s\) is the identity permutation. On the other hand, formula (1) determines \(u\cap v\) uniquely: indeed, let \((e_1',\ldots ,e_r')\) be a basis of \(\mathrm{M}'\) , \((f_1'',\dots,f_s'')\) a basis of \(\mathrm{M}''\) , and \(f_{1},\ldots ,f_{s}\) elements of \(\mathrm{M}\) such that \(p(f_i) = f_i''\) for \(1\leq i\leq s\) . Then \((f_{1},\ldots ,f_{s},i(e_1'),\ldots ,i(e_r'))\) is a basis of \(\mathrm{M}\) (Algebra, Chap. II, §1, no. 11, Prop. 21), and formula (1) can be written

\[
(u \cap v) (f _ {1}, \dots , f _ {s}, i (e _ {1} ^ {\prime}), \dots , i (e _ {r} ^ {\prime})) = u (f _ {1} ^ {\prime \prime}, \dots , f _ {s} ^ {\prime \prime}) v (e _ {1} ^ {\prime}, \dots , e _ {r} ^ {\prime});\tag{2}
\]

but an element of \(\mathrm{Alt}^{s + r}(\mathbf{M})\) is determined by its value on a basis.

It follows from the preceding that each of the conditions a) and b) determines the product \(u \cap v\) uniquely; it is clear that this product is bilinear. Finally, the last assertion of the lemma follows from formula (2).